\documentclass[11pt,a4paper]{article} \usepackage[utf8]{inputenc} \usepackage{lmodern} \usepackage[T1]{fontenc}
\usepackage[english]{babel} \usepackage[margin=2.3cm]{geometry} \usepackage{xcolor}
\usepackage{booktabs} \usepackage{longtable} \usepackage{array} \usepackage{amsmath,amssymb}
\usepackage{listings} \usepackage[hidelinks]{hyperref} \newcommand{\code}[1]{\texttt{#1}}
% Generated by tools/gen_concepts_values.py -- do not edit.
% Values come from mal-core/config/model_constants.yaml via the
% validated loader; a key that no longer exists makes \mskey fail
% the build with an error, in any document that references it.

\makeatletter
\newcommand{\mskey}[1]{%
  \expandafter\ifx\csname ms@k@#1\endcsname\relax
    \errmessage{Concept document: unknown model constant key `#1'}%
  \else
    \csname ms@k@#1\endcsname%
  \fi}
\makeatother

\expandafter\def\csname ms@k@adults.daily_mort_basal\endcsname{0.87}
\expandafter\def\csname ms@k@adults.mort_cap\endcsname{0.93}
\expandafter\def\csname ms@k@adults.mort_floor\endcsname{0.6}
\expandafter\def\csname ms@k@adults.opt_c\endcsname{25.0}
\expandafter\def\csname ms@k@adults.sigma_c\endcsname{15.0}
\expandafter\def\csname ms@k@aquatic_stages.egg_daily_mort\endcsname{0.08}
\expandafter\def\csname ms@k@aquatic_stages.larva_bh_alpha\endcsname{0.05}
\expandafter\def\csname ms@k@aquatic_stages.larva_bh_s0\endcsname{0.95}
\expandafter\def\csname ms@k@aquatic_stages.last_larval_instar\endcsname{4}
\expandafter\def\csname ms@k@aquatic_stages.pupa_daily_mort\endcsname{0.05}
\expandafter\def\csname ms@k@aquatic_stages.salinity_missing_psu\endcsname{0.0}
\expandafter\def\csname ms@k@aquatic_stress.desicc_base_daily\endcsname{0.1}
\expandafter\def\csname ms@k@aquatic_stress.desicc_grace_days\endcsname{5.0}
\expandafter\def\csname ms@k@aquatic_stress.desicc_ramp_decay\endcsname{0.3}
\expandafter\def\csname ms@k@aquatic_stress.vulnerable_instar_max\endcsname{2}
\expandafter\def\csname ms@k@aquatic_stress.wash_fraction_max\endcsname{0.6}
\expandafter\def\csname ms@k@aquatic_stress.wash_threshold_mm\endcsname{40.0}
\expandafter\def\csname ms@k@capacity.persistent_availability\endcsname{1.0}
\expandafter\def\csname ms@k@contracts.cogs_state\endcsname{"4.0"}
\expandafter\def\csname ms@k@contracts.env_nc\endcsname{"3.0"}
\expandafter\def\csname ms@k@contracts.habitat\endcsname{"4.0"}
\expandafter\def\csname ms@k@dispersal.flight_hour\endcsname{20}
\expandafter\def\csname ms@k@dispersal.night_end_hour\endcsname{6}
\expandafter\def\csname ms@k@dispersal.night_start_hour\endcsname{18}
\expandafter\def\csname ms@k@dispersal.snap_radius_m\endcsname{500.0}
\expandafter\def\csname ms@k@dispersal.wind_flight_hours\endcsname{4.0}
\expandafter\def\csname ms@k@dispersal.wind_migration_prob\endcsname{0.05}
\expandafter\def\csname ms@k@dispersal.wind_self_speed_ms\endcsname{1.0}
\expandafter\def\csname ms@k@dispersal.wind_survival\endcsname{0.85}
\expandafter\def\csname ms@k@eip.base_c\endcsname{16.0}
\expandafter\def\csname ms@k@eip.threshold_gd\endcsname{110.0}
\expandafter\def\csname ms@k@geodesy.earth_radius_m\endcsname{6371008.8}
\expandafter\def\csname ms@k@geodesy.metres_per_deg_lat\endcsname{111195.08023353292}
\expandafter\def\csname ms@k@gonotrophic.cycle_duration_days\endcsname{2.0}
\expandafter\def\csname ms@k@gonotrophic.egg_batch_max\endcsname{170}
\expandafter\def\csname ms@k@gonotrophic.egg_batch_mean\endcsname{120}
\expandafter\def\csname ms@k@gonotrophic.egg_batch_min\endcsname{30}
\expandafter\def\csname ms@k@gonotrophic.egg_batch_n\endcsname{240}
\expandafter\def\csname ms@k@gonotrophic.feeding_success_rate\endcsname{0.825}
\expandafter\def\csname ms@k@gonotrophic.first_cycle_days\endcsname{4.0}
\expandafter\def\csname ms@k@gonotrophic.resting_duration_days\endcsname{0.5}
\expandafter\def\csname ms@k@habitat.activation_min_engine\endcsname{0.01}
\expandafter\def\csname ms@k@habitat.candidate_min_engine\endcsname{0.01}
\expandafter\def\csname ms@k@habitat_types.dynamics.crop_pluvial\endcsname{"pluvial"}
\expandafter\def\csname ms@k@habitat_types.dynamics.open_pluvial\endcsname{"pluvial"}
\expandafter\def\csname ms@k@habitat_types.dynamics.surface_persistent\endcsname{"persistent"}
\expandafter\def\csname ms@k@habitat_types.dynamics.urban_persistent\endcsname{"persistent"}
\expandafter\def\csname ms@k@habitat_types.dynamics.wetland_persistent\endcsname{"persistent"}
\expandafter\def\csname ms@k@habitat_types.dynamics.wetland_pluvial\endcsname{"pluvial"}
\expandafter\def\csname ms@k@habitat_types.types\endcsname{"surface\_persistent, wetland\_persistent, urban\_persistent, wetland\_pluvial, open\_pluvial, crop\_pluvial"}
\expandafter\def\csname ms@k@host_seeking.indoor_attraction_add\endcsname{0.72}
\expandafter\def\csname ms@k@host_seeking.indoor_fraction_rural\endcsname{0.3}
\expandafter\def\csname ms@k@host_seeking.indoor_fraction_urban\endcsname{0.72}
\expandafter\def\csname ms@k@host_seeking.starvation_days\endcsname{7}
\expandafter\def\csname ms@k@host_seeking.urbanicity_add\endcsname{0.2}
\expandafter\def\csname ms@k@ingest.candidate_min_ingest\endcsname{1e-06}
\expandafter\def\csname ms@k@ingest.cell_size_hint_m\endcsname{1000.0}
\expandafter\def\csname ms@k@ingest.coastline_buffer_m\endcsname{2000.0}
\expandafter\def\csname ms@k@ingest.coastline_resolution\endcsname{"f"}
\expandafter\def\csname ms@k@ingest.fill_frac_valid_max\endcsname{0.5}
\expandafter\def\csname ms@k@ingest.grid_degrees_step_max\endcsname{1}
\expandafter\def\csname ms@k@ingest.land_mask_threshold\endcsname{0.5}
\expandafter\def\csname ms@k@ingest.min_dilation_radius_cells\endcsname{1}
\expandafter\def\csname ms@k@ingest.nc_compression_level\endcsname{4}
\expandafter\def\csname ms@k@ingest.nodata_sentinel\endcsname{-9999.0}
\expandafter\def\csname ms@k@ingest.nodata_sentinel_int\endcsname{-9999}
\expandafter\def\csname ms@k@ingest.ocean_mask_threshold\endcsname{0.5}
\expandafter\def\csname ms@k@ingest.residential_from_indoor_frac\endcsname{0.8}
\expandafter\def\csname ms@k@ingest.sea_water_floor\endcsname{0.05}
\expandafter\def\csname ms@k@ingest.valid_count_epsilon\endcsname{1e-09}
\expandafter\def\csname ms@k@integration.directed_flight_budget_h\endcsname{0.4856}
\expandafter\def\csname ms@k@integration.dt_olfactory_s\endcsname{0.25}
\expandafter\def\csname ms@k@integration.dt_terminal_s\endcsname{0.04}
\expandafter\def\csname ms@k@integration.landing_speed_hold_s\endcsname{3.0}
\expandafter\def\csname ms@k@integration.landing_speed_mm_s\endcsname{50.0}
\expandafter\def\csname ms@k@integration.phase_duration_hours\endcsname{"12.0, 4.0, 6.0, 2.0"}
\expandafter\def\csname ms@k@integration.plume_memory_s\endcsname{0.7}
\expandafter\def\csname ms@k@mating.female_attendance_rate\endcsname{0.1033}
\expandafter\def\csname ms@k@mating.local_swarm_male_offer_factor\endcsname{0.0}
\expandafter\def\csname ms@k@mating.male_females_per_night_max\endcsname{3}
\expandafter\def\csname ms@k@mating.male_move_capacity_m\endcsname{50.0}
\expandafter\def\csname ms@k@mating.male_swarm_reference_size\endcsname{100.0}
\expandafter\def\csname ms@k@mating.p_mate_male_fraction\endcsname{0.072}
\expandafter\def\csname ms@k@mating.p_mate_pairs_per_individual\endcsname{0.167}
\expandafter\def\csname ms@k@mating.p_mate_size_exponent\endcsname{0.0}
\expandafter\def\csname ms@k@mating.rain_suppresses_swarming\endcsname{true}
\expandafter\def\csname ms@k@mating.swarm_duration_min\endcsname{28.3}
\expandafter\def\csname ms@k@mating.swarm_height_m\endcsname{2.87}
\expandafter\def\csname ms@k@mating.swarm_markers_per_km2\endcsname{1.0}
\expandafter\def\csname ms@k@mating.swarm_markers_source\endcsname{"engineering prior"}
\expandafter\def\csname ms@k@mating.swarm_start_min_after_sunset\endcsname{5.0}
\expandafter\def\csname ms@k@oviposition.detect_range_m\endcsname{2000.0}
\expandafter\def\csname ms@k@oviposition.min_attraction\endcsname{0.3}
\expandafter\def\csname ms@k@oviposition.scale_m\endcsname{50.0}
\expandafter\def\csname ms@k@oviposition.step_m\endcsname{50.0}
\expandafter\def\csname ms@k@oviposition.timeout_days\endcsname{7}
\expandafter\def\csname ms@k@pluvial.cold_start_z\endcsname{0.0}
\expandafter\def\csname ms@k@pluvial.delta_days\endcsname{1.0}
\expandafter\def\csname ms@k@pluvial.dry_memory_days\endcsname{5.0}
\expandafter\def\csname ms@k@pluvial.rain_scale_mm\endcsname{2.5}
\expandafter\def\csname ms@k@rng.aquatic_seed_larva_mortality\endcsname{123}
\expandafter\def\csname ms@k@rng.aquatic_seed_stage_mortality\endcsname{42}
\expandafter\def\csname ms@k@rng.aquatic_seed_washout\endcsname{789}
\expandafter\def\csname ms@k@species.arabiensis.activity_weights\endcsname{"0.02, 0.1, 0.8, 0.08"}
\expandafter\def\csname ms@k@species.arabiensis.aestivation_supported\endcsname{true}
\expandafter\def\csname ms@k@species.arabiensis.anemotaxis_deviation_deg\endcsname{30.0}
\expandafter\def\csname ms@k@species.arabiensis.anemotaxis_rule\endcsname{"upwind"}
\expandafter\def\csname ms@k@species.arabiensis.arrested_daily_survival\endcsname{0.986}
\expandafter\def\csname ms@k@species.arabiensis.blood_meal_refractory_days\endcsname{0.5}
\expandafter\def\csname ms@k@species.arabiensis.blood_meal_size_ul\endcsname{0.3}
\expandafter\def\csname ms@k@species.arabiensis.blood_meal_target_ul\endcsname{0.5}
\expandafter\def\csname ms@k@species.arabiensis.casting_angle_deg\endcsname{90.0}
\expandafter\def\csname ms@k@species.arabiensis.casting_probability\endcsname{0.1}
\expandafter\def\csname ms@k@species.arabiensis.co2_activation_radius_m\endcsname{100.0}
\expandafter\def\csname ms@k@species.arabiensis.daily_move_capacity_m\endcsname{590.0}
\expandafter\def\csname ms@k@species.arabiensis.directed_flight_budget_h\endcsname{0.4856}
\expandafter\def\csname ms@k@species.arabiensis.eggs_per_ul\endcsname{223.0}
\expandafter\def\csname ms@k@species.arabiensis.exploration_range_m\endcsname{100.0}
\expandafter\def\csname ms@k@species.arabiensis.exploration_step_correlation\endcsname{0.5}
\expandafter\def\csname ms@k@species.arabiensis.exploration_step_m\endcsname{50.0}
\expandafter\def\csname ms@k@species.arabiensis.first_cycle_target_ul\endcsname{0.8}
\expandafter\def\csname ms@k@species.arabiensis.flight_speed_ms\endcsname{0.3375}
\expandafter\def\csname ms@k@species.arabiensis.gonotrophic_arrest_prob_dry\endcsname{0.4454}
\expandafter\def\csname ms@k@species.arabiensis.male_daily_survival\endcsname{0.773}
\expandafter\def\csname ms@k@species.arabiensis.male_ratio\endcsname{0.5}
\expandafter\def\csname ms@k@species.arabiensis.meal_size_min_ul\endcsname{0.1}
\expandafter\def\csname ms@k@species.arabiensis.olfactory_radius_m\endcsname{50.0}
\expandafter\def\csname ms@k@species.arabiensis.plume_memory_s\endcsname{0.7}
\expandafter\def\csname ms@k@species.arabiensis.pref.cattle\endcsname{0.2}
\expandafter\def\csname ms@k@species.arabiensis.pref.goat\endcsname{0.15}
\expandafter\def\csname ms@k@species.arabiensis.pref.human\endcsname{0.5}
\expandafter\def\csname ms@k@species.arabiensis.pref.sheep\endcsname{0.1}
\expandafter\def\csname ms@k@species.arabiensis.pref.wildlife\endcsname{0.05}
\expandafter\def\csname ms@k@species.arabiensis.probing_success_rate\endcsname{0.7275}
\expandafter\def\csname ms@k@species.arabiensis.salinity_hi_tol_ppt\endcsname{15.0}
\expandafter\def\csname ms@k@species.arabiensis.salinity_opt_ppt\endcsname{0.0}
\expandafter\def\csname ms@k@species.arabiensis.salinity_slope\endcsname{0.15}
\expandafter\def\csname ms@k@species.arabiensis.terminal_radius_m\endcsname{0.15}
\expandafter\def\csname ms@k@species.arabiensis.visual_radius_m\endcsname{0.485}
\expandafter\def\csname ms@k@species.arabiensis.weights.crop_pluvial\endcsname{0.1}
\expandafter\def\csname ms@k@species.arabiensis.weights.open_pluvial\endcsname{0.55}
\expandafter\def\csname ms@k@species.arabiensis.weights.surface_persistent\endcsname{0.1}
\expandafter\def\csname ms@k@species.arabiensis.weights.urban_persistent\endcsname{0.15}
\expandafter\def\csname ms@k@species.arabiensis.weights.wetland_persistent\endcsname{0.1}
\expandafter\def\csname ms@k@species.arabiensis.weights.wetland_pluvial\endcsname{0.0}
\expandafter\def\csname ms@k@species.arabiensis.wind_block_threshold_ms\endcsname{0.3375}
\expandafter\def\csname ms@k@species.arabiensis.wind_reference\endcsname{"u100/v100"}
\expandafter\def\csname ms@k@species.coluzzii.activity_weights\endcsname{"0.0, 0.255, 0.48, 0.265"}
\expandafter\def\csname ms@k@species.coluzzii.aestivation_supported\endcsname{true}
\expandafter\def\csname ms@k@species.coluzzii.anemotaxis_deviation_deg\endcsname{30.0}
\expandafter\def\csname ms@k@species.coluzzii.anemotaxis_rule\endcsname{"upwind"}
\expandafter\def\csname ms@k@species.coluzzii.arrested_daily_survival\endcsname{0.986}
\expandafter\def\csname ms@k@species.coluzzii.blood_meal_refractory_days\endcsname{0.5}
\expandafter\def\csname ms@k@species.coluzzii.blood_meal_size_ul\endcsname{0.3}
\expandafter\def\csname ms@k@species.coluzzii.blood_meal_target_ul\endcsname{0.5}
\expandafter\def\csname ms@k@species.coluzzii.casting_angle_deg\endcsname{90.0}
\expandafter\def\csname ms@k@species.coluzzii.casting_probability\endcsname{0.1}
\expandafter\def\csname ms@k@species.coluzzii.co2_activation_radius_m\endcsname{100.0}
\expandafter\def\csname ms@k@species.coluzzii.daily_move_capacity_m\endcsname{590.0}
\expandafter\def\csname ms@k@species.coluzzii.directed_flight_budget_h\endcsname{0.4856}
\expandafter\def\csname ms@k@species.coluzzii.eggs_per_ul\endcsname{223.0}
\expandafter\def\csname ms@k@species.coluzzii.exploration_range_m\endcsname{100.0}
\expandafter\def\csname ms@k@species.coluzzii.exploration_step_correlation\endcsname{0.5}
\expandafter\def\csname ms@k@species.coluzzii.exploration_step_m\endcsname{50.0}
\expandafter\def\csname ms@k@species.coluzzii.first_cycle_target_ul\endcsname{0.8}
\expandafter\def\csname ms@k@species.coluzzii.flight_speed_ms\endcsname{0.3375}
\expandafter\def\csname ms@k@species.coluzzii.gonotrophic_arrest_prob_dry\endcsname{0.2704}
\expandafter\def\csname ms@k@species.coluzzii.male_daily_survival\endcsname{0.773}
\expandafter\def\csname ms@k@species.coluzzii.male_ratio\endcsname{0.5}
\expandafter\def\csname ms@k@species.coluzzii.meal_size_min_ul\endcsname{0.1}
\expandafter\def\csname ms@k@species.coluzzii.olfactory_radius_m\endcsname{50.0}
\expandafter\def\csname ms@k@species.coluzzii.plume_memory_s\endcsname{0.7}
\expandafter\def\csname ms@k@species.coluzzii.pref.cattle\endcsname{0.005}
\expandafter\def\csname ms@k@species.coluzzii.pref.goat\endcsname{0.003}
\expandafter\def\csname ms@k@species.coluzzii.pref.human\endcsname{0.99}
\expandafter\def\csname ms@k@species.coluzzii.pref.sheep\endcsname{0.001}
\expandafter\def\csname ms@k@species.coluzzii.pref.wildlife\endcsname{0.001}
\expandafter\def\csname ms@k@species.coluzzii.probing_success_rate\endcsname{0.7275}
\expandafter\def\csname ms@k@species.coluzzii.salinity_hi_tol_ppt\endcsname{30.0}
\expandafter\def\csname ms@k@species.coluzzii.salinity_opt_ppt\endcsname{0.0}
\expandafter\def\csname ms@k@species.coluzzii.salinity_slope\endcsname{0.08}
\expandafter\def\csname ms@k@species.coluzzii.terminal_radius_m\endcsname{0.15}
\expandafter\def\csname ms@k@species.coluzzii.visual_radius_m\endcsname{0.485}
\expandafter\def\csname ms@k@species.coluzzii.weights.crop_pluvial\endcsname{0.05}
\expandafter\def\csname ms@k@species.coluzzii.weights.open_pluvial\endcsname{0.25}
\expandafter\def\csname ms@k@species.coluzzii.weights.surface_persistent\endcsname{0.4}
\expandafter\def\csname ms@k@species.coluzzii.weights.urban_persistent\endcsname{0.2}
\expandafter\def\csname ms@k@species.coluzzii.weights.wetland_persistent\endcsname{0.1}
\expandafter\def\csname ms@k@species.coluzzii.weights.wetland_pluvial\endcsname{0.0}
\expandafter\def\csname ms@k@species.coluzzii.wind_block_threshold_ms\endcsname{0.3375}
\expandafter\def\csname ms@k@species.coluzzii.wind_reference\endcsname{"u100/v100"}
\expandafter\def\csname ms@k@species.default_species\endcsname{"coluzzii"}
\expandafter\def\csname ms@k@species.funestus.activity_weights\endcsname{"0.02, 0.1, 0.8, 0.08"}
\expandafter\def\csname ms@k@species.funestus.aestivation_supported\endcsname{false}
\expandafter\def\csname ms@k@species.funestus.anemotaxis_deviation_deg\endcsname{30.0}
\expandafter\def\csname ms@k@species.funestus.anemotaxis_rule\endcsname{"upwind"}
\expandafter\def\csname ms@k@species.funestus.arrested_daily_survival\endcsname{0.986}
\expandafter\def\csname ms@k@species.funestus.blood_meal_refractory_days\endcsname{0.5}
\expandafter\def\csname ms@k@species.funestus.blood_meal_size_ul\endcsname{0.3}
\expandafter\def\csname ms@k@species.funestus.blood_meal_target_ul\endcsname{0.5}
\expandafter\def\csname ms@k@species.funestus.casting_angle_deg\endcsname{90.0}
\expandafter\def\csname ms@k@species.funestus.casting_probability\endcsname{0.1}
\expandafter\def\csname ms@k@species.funestus.co2_activation_radius_m\endcsname{100.0}
\expandafter\def\csname ms@k@species.funestus.daily_move_capacity_m\endcsname{590.0}
\expandafter\def\csname ms@k@species.funestus.directed_flight_budget_h\endcsname{0.4856}
\expandafter\def\csname ms@k@species.funestus.eggs_per_ul\endcsname{223.0}
\expandafter\def\csname ms@k@species.funestus.exploration_range_m\endcsname{100.0}
\expandafter\def\csname ms@k@species.funestus.exploration_step_correlation\endcsname{0.5}
\expandafter\def\csname ms@k@species.funestus.exploration_step_m\endcsname{50.0}
\expandafter\def\csname ms@k@species.funestus.first_cycle_target_ul\endcsname{0.8}
\expandafter\def\csname ms@k@species.funestus.flight_speed_ms\endcsname{0.3375}
\expandafter\def\csname ms@k@species.funestus.gonotrophic_arrest_prob_dry\endcsname{0.0}
\expandafter\def\csname ms@k@species.funestus.male_daily_survival\endcsname{0.773}
\expandafter\def\csname ms@k@species.funestus.male_ratio\endcsname{0.5}
\expandafter\def\csname ms@k@species.funestus.meal_size_min_ul\endcsname{0.1}
\expandafter\def\csname ms@k@species.funestus.olfactory_radius_m\endcsname{50.0}
\expandafter\def\csname ms@k@species.funestus.plume_memory_s\endcsname{0.7}
\expandafter\def\csname ms@k@species.funestus.pref.cattle\endcsname{0.005}
\expandafter\def\csname ms@k@species.funestus.pref.goat\endcsname{0.003}
\expandafter\def\csname ms@k@species.funestus.pref.human\endcsname{0.99}
\expandafter\def\csname ms@k@species.funestus.pref.sheep\endcsname{0.001}
\expandafter\def\csname ms@k@species.funestus.pref.wildlife\endcsname{0.001}
\expandafter\def\csname ms@k@species.funestus.probing_success_rate\endcsname{0.7275}
\expandafter\def\csname ms@k@species.funestus.salinity_hi_tol_ppt\endcsname{15.0}
\expandafter\def\csname ms@k@species.funestus.salinity_opt_ppt\endcsname{0.0}
\expandafter\def\csname ms@k@species.funestus.salinity_slope\endcsname{0.15}
\expandafter\def\csname ms@k@species.funestus.terminal_radius_m\endcsname{0.15}
\expandafter\def\csname ms@k@species.funestus.visual_radius_m\endcsname{0.485}
\expandafter\def\csname ms@k@species.funestus.weights.crop_pluvial\endcsname{0.05}
\expandafter\def\csname ms@k@species.funestus.weights.open_pluvial\endcsname{0.1}
\expandafter\def\csname ms@k@species.funestus.weights.surface_persistent\endcsname{0.45}
\expandafter\def\csname ms@k@species.funestus.weights.urban_persistent\endcsname{0.1}
\expandafter\def\csname ms@k@species.funestus.weights.wetland_persistent\endcsname{0.3}
\expandafter\def\csname ms@k@species.funestus.weights.wetland_pluvial\endcsname{0.0}
\expandafter\def\csname ms@k@species.funestus.wind_block_threshold_ms\endcsname{0.3375}
\expandafter\def\csname ms@k@species.funestus.wind_reference\endcsname{"u100/v100"}
\expandafter\def\csname ms@k@species.gambiae_ss.activity_weights\endcsname{"0.02, 0.1, 0.8, 0.08"}
\expandafter\def\csname ms@k@species.gambiae_ss.aestivation_supported\endcsname{true}
\expandafter\def\csname ms@k@species.gambiae_ss.anemotaxis_deviation_deg\endcsname{30.0}
\expandafter\def\csname ms@k@species.gambiae_ss.anemotaxis_rule\endcsname{"upwind"}
\expandafter\def\csname ms@k@species.gambiae_ss.arrested_daily_survival\endcsname{0.986}
\expandafter\def\csname ms@k@species.gambiae_ss.blood_meal_refractory_days\endcsname{0.5}
\expandafter\def\csname ms@k@species.gambiae_ss.blood_meal_size_ul\endcsname{0.3}
\expandafter\def\csname ms@k@species.gambiae_ss.blood_meal_target_ul\endcsname{0.5}
\expandafter\def\csname ms@k@species.gambiae_ss.casting_angle_deg\endcsname{90.0}
\expandafter\def\csname ms@k@species.gambiae_ss.casting_probability\endcsname{0.1}
\expandafter\def\csname ms@k@species.gambiae_ss.co2_activation_radius_m\endcsname{100.0}
\expandafter\def\csname ms@k@species.gambiae_ss.daily_move_capacity_m\endcsname{590.0}
\expandafter\def\csname ms@k@species.gambiae_ss.directed_flight_budget_h\endcsname{0.4856}
\expandafter\def\csname ms@k@species.gambiae_ss.eggs_per_ul\endcsname{223.0}
\expandafter\def\csname ms@k@species.gambiae_ss.exploration_range_m\endcsname{100.0}
\expandafter\def\csname ms@k@species.gambiae_ss.exploration_step_correlation\endcsname{0.5}
\expandafter\def\csname ms@k@species.gambiae_ss.exploration_step_m\endcsname{50.0}
\expandafter\def\csname ms@k@species.gambiae_ss.first_cycle_target_ul\endcsname{0.8}
\expandafter\def\csname ms@k@species.gambiae_ss.flight_speed_ms\endcsname{0.3375}
\expandafter\def\csname ms@k@species.gambiae_ss.gonotrophic_arrest_prob_dry\endcsname{0.6838}
\expandafter\def\csname ms@k@species.gambiae_ss.male_daily_survival\endcsname{0.773}
\expandafter\def\csname ms@k@species.gambiae_ss.male_ratio\endcsname{0.5}
\expandafter\def\csname ms@k@species.gambiae_ss.meal_size_min_ul\endcsname{0.1}
\expandafter\def\csname ms@k@species.gambiae_ss.olfactory_radius_m\endcsname{50.0}
\expandafter\def\csname ms@k@species.gambiae_ss.plume_memory_s\endcsname{0.7}
\expandafter\def\csname ms@k@species.gambiae_ss.pref.cattle\endcsname{0.005}
\expandafter\def\csname ms@k@species.gambiae_ss.pref.goat\endcsname{0.003}
\expandafter\def\csname ms@k@species.gambiae_ss.pref.human\endcsname{0.99}
\expandafter\def\csname ms@k@species.gambiae_ss.pref.sheep\endcsname{0.001}
\expandafter\def\csname ms@k@species.gambiae_ss.pref.wildlife\endcsname{0.001}
\expandafter\def\csname ms@k@species.gambiae_ss.probing_success_rate\endcsname{0.7275}
\expandafter\def\csname ms@k@species.gambiae_ss.salinity_hi_tol_ppt\endcsname{20.0}
\expandafter\def\csname ms@k@species.gambiae_ss.salinity_opt_ppt\endcsname{0.0}
\expandafter\def\csname ms@k@species.gambiae_ss.salinity_slope\endcsname{0.12}
\expandafter\def\csname ms@k@species.gambiae_ss.terminal_radius_m\endcsname{0.15}
\expandafter\def\csname ms@k@species.gambiae_ss.visual_radius_m\endcsname{0.485}
\expandafter\def\csname ms@k@species.gambiae_ss.weights.crop_pluvial\endcsname{0.1}
\expandafter\def\csname ms@k@species.gambiae_ss.weights.open_pluvial\endcsname{0.65}
\expandafter\def\csname ms@k@species.gambiae_ss.weights.surface_persistent\endcsname{0.05}
\expandafter\def\csname ms@k@species.gambiae_ss.weights.urban_persistent\endcsname{0.15}
\expandafter\def\csname ms@k@species.gambiae_ss.weights.wetland_persistent\endcsname{0.05}
\expandafter\def\csname ms@k@species.gambiae_ss.weights.wetland_pluvial\endcsname{0.0}
\expandafter\def\csname ms@k@species.gambiae_ss.wind_block_threshold_ms\endcsname{0.3375}
\expandafter\def\csname ms@k@species.gambiae_ss.wind_reference\endcsname{"u100/v100"}
\expandafter\def\csname ms@k@species.k_ref\endcsname{50.0}
\expandafter\def\csname ms@k@species.melas.activity_weights\endcsname{"0.02, 0.1, 0.8, 0.08"}
\expandafter\def\csname ms@k@species.melas.aestivation_supported\endcsname{false}
\expandafter\def\csname ms@k@species.melas.anemotaxis_deviation_deg\endcsname{30.0}
\expandafter\def\csname ms@k@species.melas.anemotaxis_rule\endcsname{"upwind"}
\expandafter\def\csname ms@k@species.melas.arrested_daily_survival\endcsname{0.986}
\expandafter\def\csname ms@k@species.melas.blood_meal_refractory_days\endcsname{0.5}
\expandafter\def\csname ms@k@species.melas.blood_meal_size_ul\endcsname{0.3}
\expandafter\def\csname ms@k@species.melas.blood_meal_target_ul\endcsname{0.5}
\expandafter\def\csname ms@k@species.melas.casting_angle_deg\endcsname{90.0}
\expandafter\def\csname ms@k@species.melas.casting_probability\endcsname{0.1}
\expandafter\def\csname ms@k@species.melas.co2_activation_radius_m\endcsname{100.0}
\expandafter\def\csname ms@k@species.melas.daily_move_capacity_m\endcsname{590.0}
\expandafter\def\csname ms@k@species.melas.directed_flight_budget_h\endcsname{0.4856}
\expandafter\def\csname ms@k@species.melas.eggs_per_ul\endcsname{223.0}
\expandafter\def\csname ms@k@species.melas.exploration_range_m\endcsname{100.0}
\expandafter\def\csname ms@k@species.melas.exploration_step_correlation\endcsname{0.5}
\expandafter\def\csname ms@k@species.melas.exploration_step_m\endcsname{50.0}
\expandafter\def\csname ms@k@species.melas.first_cycle_target_ul\endcsname{0.8}
\expandafter\def\csname ms@k@species.melas.flight_speed_ms\endcsname{0.3375}
\expandafter\def\csname ms@k@species.melas.gonotrophic_arrest_prob_dry\endcsname{0.0}
\expandafter\def\csname ms@k@species.melas.male_daily_survival\endcsname{0.773}
\expandafter\def\csname ms@k@species.melas.male_ratio\endcsname{0.5}
\expandafter\def\csname ms@k@species.melas.meal_size_min_ul\endcsname{0.1}
\expandafter\def\csname ms@k@species.melas.olfactory_radius_m\endcsname{50.0}
\expandafter\def\csname ms@k@species.melas.plume_memory_s\endcsname{0.7}
\expandafter\def\csname ms@k@species.melas.pref.cattle\endcsname{0.005}
\expandafter\def\csname ms@k@species.melas.pref.goat\endcsname{0.003}
\expandafter\def\csname ms@k@species.melas.pref.human\endcsname{0.99}
\expandafter\def\csname ms@k@species.melas.pref.sheep\endcsname{0.001}
\expandafter\def\csname ms@k@species.melas.pref.wildlife\endcsname{0.001}
\expandafter\def\csname ms@k@species.melas.probing_success_rate\endcsname{0.7275}
\expandafter\def\csname ms@k@species.melas.salinity_hi_tol_ppt\endcsname{38.0}
\expandafter\def\csname ms@k@species.melas.salinity_opt_ppt\endcsname{0.0}
\expandafter\def\csname ms@k@species.melas.salinity_slope\endcsname{0.02}
\expandafter\def\csname ms@k@species.melas.terminal_radius_m\endcsname{0.15}
\expandafter\def\csname ms@k@species.melas.visual_radius_m\endcsname{0.485}
\expandafter\def\csname ms@k@species.melas.weights.crop_pluvial\endcsname{0.05}
\expandafter\def\csname ms@k@species.melas.weights.open_pluvial\endcsname{0.4}
\expandafter\def\csname ms@k@species.melas.weights.surface_persistent\endcsname{0.25}
\expandafter\def\csname ms@k@species.melas.weights.urban_persistent\endcsname{0.15}
\expandafter\def\csname ms@k@species.melas.weights.wetland_persistent\endcsname{0.15}
\expandafter\def\csname ms@k@species.melas.weights.wetland_pluvial\endcsname{0.0}
\expandafter\def\csname ms@k@species.melas.wind_block_threshold_ms\endcsname{0.3375}
\expandafter\def\csname ms@k@species.melas.wind_reference\endcsname{"u100/v100"}
\expandafter\def\csname ms@k@species.neutral.weights.crop_pluvial\endcsname{0.05}
\expandafter\def\csname ms@k@species.neutral.weights.open_pluvial\endcsname{0.4}
\expandafter\def\csname ms@k@species.neutral.weights.surface_persistent\endcsname{0.25}
\expandafter\def\csname ms@k@species.neutral.weights.urban_persistent\endcsname{0.15}
\expandafter\def\csname ms@k@species.neutral.weights.wetland_persistent\endcsname{0.15}
\expandafter\def\csname ms@k@species.neutral.weights.wetland_pluvial\endcsname{0.0}
\expandafter\def\csname ms@k@species.stephensi.activity_weights\endcsname{"0.02, 0.1, 0.8, 0.08"}
\expandafter\def\csname ms@k@species.stephensi.aestivation_supported\endcsname{false}
\expandafter\def\csname ms@k@species.stephensi.anemotaxis_deviation_deg\endcsname{30.0}
\expandafter\def\csname ms@k@species.stephensi.anemotaxis_rule\endcsname{"upwind"}
\expandafter\def\csname ms@k@species.stephensi.arrested_daily_survival\endcsname{0.986}
\expandafter\def\csname ms@k@species.stephensi.blood_meal_refractory_days\endcsname{0.5}
\expandafter\def\csname ms@k@species.stephensi.blood_meal_size_ul\endcsname{0.3}
\expandafter\def\csname ms@k@species.stephensi.blood_meal_target_ul\endcsname{0.5}
\expandafter\def\csname ms@k@species.stephensi.casting_angle_deg\endcsname{90.0}
\expandafter\def\csname ms@k@species.stephensi.casting_probability\endcsname{0.1}
\expandafter\def\csname ms@k@species.stephensi.co2_activation_radius_m\endcsname{100.0}
\expandafter\def\csname ms@k@species.stephensi.daily_move_capacity_m\endcsname{590.0}
\expandafter\def\csname ms@k@species.stephensi.directed_flight_budget_h\endcsname{0.4856}
\expandafter\def\csname ms@k@species.stephensi.eggs_per_ul\endcsname{223.0}
\expandafter\def\csname ms@k@species.stephensi.exploration_range_m\endcsname{100.0}
\expandafter\def\csname ms@k@species.stephensi.exploration_step_correlation\endcsname{0.5}
\expandafter\def\csname ms@k@species.stephensi.exploration_step_m\endcsname{50.0}
\expandafter\def\csname ms@k@species.stephensi.first_cycle_target_ul\endcsname{0.8}
\expandafter\def\csname ms@k@species.stephensi.flight_speed_ms\endcsname{0.3375}
\expandafter\def\csname ms@k@species.stephensi.gonotrophic_arrest_prob_dry\endcsname{0.0}
\expandafter\def\csname ms@k@species.stephensi.male_daily_survival\endcsname{0.773}
\expandafter\def\csname ms@k@species.stephensi.male_ratio\endcsname{0.5}
\expandafter\def\csname ms@k@species.stephensi.meal_size_min_ul\endcsname{0.1}
\expandafter\def\csname ms@k@species.stephensi.olfactory_radius_m\endcsname{50.0}
\expandafter\def\csname ms@k@species.stephensi.plume_memory_s\endcsname{0.7}
\expandafter\def\csname ms@k@species.stephensi.pref.cattle\endcsname{0.15}
\expandafter\def\csname ms@k@species.stephensi.pref.goat\endcsname{0.05}
\expandafter\def\csname ms@k@species.stephensi.pref.human\endcsname{0.7}
\expandafter\def\csname ms@k@species.stephensi.pref.sheep\endcsname{0.05}
\expandafter\def\csname ms@k@species.stephensi.pref.wildlife\endcsname{0.05}
\expandafter\def\csname ms@k@species.stephensi.probing_success_rate\endcsname{0.7275}
\expandafter\def\csname ms@k@species.stephensi.salinity_hi_tol_ppt\endcsname{20.0}
\expandafter\def\csname ms@k@species.stephensi.salinity_opt_ppt\endcsname{0.0}
\expandafter\def\csname ms@k@species.stephensi.salinity_slope\endcsname{0.15}
\expandafter\def\csname ms@k@species.stephensi.terminal_radius_m\endcsname{0.15}
\expandafter\def\csname ms@k@species.stephensi.visual_radius_m\endcsname{0.485}
\expandafter\def\csname ms@k@species.stephensi.weights.crop_pluvial\endcsname{0.05}
\expandafter\def\csname ms@k@species.stephensi.weights.open_pluvial\endcsname{0.4}
\expandafter\def\csname ms@k@species.stephensi.weights.surface_persistent\endcsname{0.25}
\expandafter\def\csname ms@k@species.stephensi.weights.urban_persistent\endcsname{0.15}
\expandafter\def\csname ms@k@species.stephensi.weights.wetland_persistent\endcsname{0.15}
\expandafter\def\csname ms@k@species.stephensi.weights.wetland_pluvial\endcsname{0.0}
\expandafter\def\csname ms@k@species.stephensi.wind_block_threshold_ms\endcsname{0.3375}
\expandafter\def\csname ms@k@species.stephensi.wind_reference\endcsname{"u100/v100"}
\expandafter\def\csname ms@k@thermal.efd_max_eggs_per_day\endcsname{0.5}
\expandafter\def\csname ms@k@thermal.efd_tmax_c\endcsname{34.0}
\expandafter\def\csname ms@k@thermal.efd_tmin_c\endcsname{16.0}
\expandafter\def\csname ms@k@thermal.efd_topt_c\endcsname{25.0}
\expandafter\def\csname ms@k@thermal.egg_development_scale\endcsname{0.00062}
\expandafter\def\csname ms@k@thermal.egg_development_tmax_c\endcsname{40.0}
\expandafter\def\csname ms@k@thermal.egg_development_tmin_c\endcsname{10.0}
\expandafter\def\csname ms@k@thermal.larva_development_scale\endcsname{0.001}
\expandafter\def\csname ms@k@thermal.larva_development_tmax_c\endcsname{40.0}
\expandafter\def\csname ms@k@thermal.larva_development_tmin_c\endcsname{10.0}
\expandafter\def\csname ms@k@thermal.pupa_development_scale\endcsname{0.00066}
\expandafter\def\csname ms@k@thermal.pupa_development_tmax_c\endcsname{40.0}
\expandafter\def\csname ms@k@thermal.pupa_development_tmin_c\endcsname{12.0}
\expandafter\def\csname ms@k@transmission.beta_hv\endcsname{0.4}
\expandafter\def\csname ms@k@transmission.beta_vh\endcsname{0.5}
\expandafter\def\csname ms@k@transmission.human_immunity_days\endcsname{180}
\expandafter\def\csname ms@k@transmission.human_incubation_days\endcsname{12}
\expandafter\def\csname ms@k@transmission.human_infectious_days\endcsname{20}
\expandafter\def\csname ms@k@transmission.initial_human_prevalence\endcsname{0.05}
\expandafter\def\csname ms@k@transmission.initial_vector_infected_frac\endcsname{0.0}
\expandafter\def\csname ms@k@urban.beta_urban\endcsname{0.75}
\expandafter\def\csname ms@k@urban.rain_pools_built_damage\endcsname{0.3}
\expandafter\def\csname ms@k@version\endcsname{3}
\expandafter\def\csname ms@k@wildlife.class_suitability.bare\endcsname{0.2}
\expandafter\def\csname ms@k@wildlife.class_suitability.built\endcsname{0.05}
\expandafter\def\csname ms@k@wildlife.class_suitability.crops\endcsname{0.3}
\expandafter\def\csname ms@k@wildlife.class_suitability.flooded_vegetation\endcsname{0.45}
\expandafter\def\csname ms@k@wildlife.class_suitability.grass\endcsname{0.6}
\expandafter\def\csname ms@k@wildlife.class_suitability.shrub_and_scrub\endcsname{0.7}
\expandafter\def\csname ms@k@wildlife.class_suitability.snow_and_ice\endcsname{0.1}
\expandafter\def\csname ms@k@wildlife.class_suitability.trees\endcsname{0.9}
\expandafter\def\csname ms@k@wildlife.class_suitability.water\endcsname{0.0}
\expandafter\def\csname ms@k@wildlife.weights.habitat\endcsname{0.5}
\expandafter\def\csname ms@k@wildlife.weights.remoteness\endcsname{0.2}
\expandafter\def\csname ms@k@wildlife.weights.water\endcsname{0.3}

% Generated by tools/gen_concepts_values.py -- do not edit.
% Values are read from the CODE that implements each stage, by AST.
% A key that no longer exists makes \codeconst fail the LaTeX build
% with an error, in any document that references it.
% Each line carries the file:line the value was read from.

\makeatletter
\newcommand{\codeconst}[1]{%
  \expandafter\ifx\csname cc@k@#1\endcsname\relax
    \errmessage{Concept document: unknown code constant `#1'}%
  \else
    \csname cc@k@#1\endcsname%
  \fi}
\makeatother

% download.wildlife_habitat_fallback = 0.5  (mal-commonlib/src/mal_commonlib/data/loaders/wildlife.py:205)
\expandafter\def\csname cc@k@download.wildlife_habitat_fallback\endcsname{0.5}
% download.wildlife_remoteness_neutral = 1.0  (mal-commonlib/src/mal_commonlib/data/loaders/wildlife.py:161)
\expandafter\def\csname cc@k@download.wildlife_remoteness_neutral\endcsname{1.0}
% ingest.bbox_east = 1.5  (mal-commonlib/src/mal_commonlib/aois.yaml:18)
\expandafter\def\csname cc@k@ingest.bbox_east\endcsname{1.5}
% ingest.bbox_north = 11.5  (mal-commonlib/src/mal_commonlib/aois.yaml:18)
\expandafter\def\csname cc@k@ingest.bbox_north\endcsname{11.5}
% ingest.bbox_south = 4.5  (mal-commonlib/src/mal_commonlib/aois.yaml:18)
\expandafter\def\csname cc@k@ingest.bbox_south\endcsname{4.5}
% ingest.bbox_west = -3.5  (mal-commonlib/src/mal_commonlib/aois.yaml:18)
\expandafter\def\csname cc@k@ingest.bbox_west\endcsname{-3.5}
% ingest.cell_size_m = 1000  (mal-commonlib/src/mal_commonlib/aois.yaml:20)
\expandafter\def\csname cc@k@ingest.cell_size_m\endcsname{1000}
% ingest.env_daily_channels = 2  (mal-core/src/mal_core/ingest/daily_nc.py:414)
\expandafter\def\csname cc@k@ingest.env_daily_channels\endcsname{2}
% ingest.env_static_planes = 5  (mal-core/src/mal_core/ingest/daily_nc.py:414)
\expandafter\def\csname cc@k@ingest.env_static_planes\endcsname{5}
% ingest.fallback_building_rural = 0.1  (mal-core/src/mal_core/ingest/hosts.py:142)
\expandafter\def\csname cc@k@ingest.fallback_building_rural\endcsname{0.1}
% ingest.fallback_building_urban = 0.6  (mal-core/src/mal_core/ingest/hosts.py:142)
\expandafter\def\csname cc@k@ingest.fallback_building_urban\endcsname{0.6}
% ingest.fallback_wildlife_proxy = 0.3  (mal-core/src/mal_core/ingest/hosts.py:165)
\expandafter\def\csname cc@k@ingest.fallback_wildlife_proxy\endcsname{0.3}
% ingest.jrc_occurrence_threshold_pct = 80  (mal-commonlib/src/mal_commonlib/data/loaders/jrc_gsw.py:72)
\expandafter\def\csname cc@k@ingest.jrc_occurrence_threshold_pct\endcsname{80}
% ingest.k_ref_diag_per_km2 = 1329900.0  (mal-core/src/mal_core/habitat/capacity.py:30)
\expandafter\def\csname cc@k@ingest.k_ref_diag_per_km2\endcsname{1329900.0}
% ingest.mobility_beta_day = 0.05  (mal-core/src/mal_core/ingest/mobility.py:43)
\expandafter\def\csname cc@k@ingest.mobility_beta_day\endcsname{0.05}
% ingest.mobility_beta_livestock = 0.1  (mal-core/src/mal_core/ingest/mobility.py:45)
\expandafter\def\csname cc@k@ingest.mobility_beta_livestock\endcsname{0.1}
% ingest.mobility_beta_night = 0.5  (mal-core/src/mal_core/ingest/mobility.py:44)
\expandafter\def\csname cc@k@ingest.mobility_beta_night\endcsname{0.5}
% ingest.mobility_cell_size_km = 1.0  (mal-core/src/mal_core/ingest/mobility.py:42)
\expandafter\def\csname cc@k@ingest.mobility_cell_size_km\endcsname{1.0}
% ingest.mobility_max_distance_km = 50.0  (mal-core/src/mal_core/ingest/mobility.py:46)
\expandafter\def\csname cc@k@ingest.mobility_max_distance_km\endcsname{50.0}
% ingest.mobility_sparsity_threshold = 0.0005  (mal-core/src/mal_core/ingest/mobility.py:47)
\expandafter\def\csname cc@k@ingest.mobility_sparsity_threshold\endcsname{0.0005}
% ingest.urban_class_rural = 50  (mal-core/src/mal_core/ingest/hosts.py:126)
\expandafter\def\csname cc@k@ingest.urban_class_rural\endcsname{50}
% ingest.urban_class_urban = 30  (mal-core/src/mal_core/ingest/hosts.py:125)
\expandafter\def\csname cc@k@ingest.urban_class_urban\endcsname{30}
% ingest.urban_class_water = 20  (mal-core/src/mal_core/ingest/hosts.py:127)
\expandafter\def\csname cc@k@ingest.urban_class_water\endcsname{20}
% prediction.block_size = 128  (mal-core/src/mal_core/prediction/predictor.py:126)
\expandafter\def\csname cc@k@prediction.block_size\endcsname{128}
% prediction.contract_version_manifest = "1.1"  (mal-core/src/mal_core/prediction/registry.py:31)
\expandafter\def\csname cc@k@prediction.contract_version_manifest\endcsname{"1.1"}
% prediction.contract_version_sidecar = "1.1"  (mal-core/src/mal_core/prediction/predictor.py:146)
\expandafter\def\csname cc@k@prediction.contract_version_sidecar\endcsname{"1.1"}
% prediction.in_channels = 10  (mal-core/src/mal_core/prediction/registry.py:32)
\expandafter\def\csname cc@k@prediction.in_channels\endcsname{10}
% prediction.manifest_version = "1.0"  (mal-core/src/mal_core/prediction/registry.py:30)
\expandafter\def\csname cc@k@prediction.manifest_version\endcsname{"1.0"}
% prediction.out_channels = 6  (mal-core/src/mal_core/prediction/registry.py:33)
\expandafter\def\csname cc@k@prediction.out_channels\endcsname{6}
% prediction.scale_continental_m = 5000  (mal-core/src/mal_core/prediction/aggregator.py:29)
\expandafter\def\csname cc@k@prediction.scale_continental_m\endcsname{5000}
% prediction.scale_national_m = 2000  (mal-core/src/mal_core/prediction/aggregator.py:28)
\expandafter\def\csname cc@k@prediction.scale_national_m\endcsname{2000}
% prediction.scale_regional_m = 1000  (mal-core/src/mal_core/prediction/aggregator.py:27)
\expandafter\def\csname cc@k@prediction.scale_regional_m\endcsname{1000}
% scoring.d11_larva_band_max = 0.9  (mal-core/src/mal_core/scoring/scorers/d11_aquatic_structure.py:35)
\expandafter\def\csname cc@k@scoring.d11_larva_band_max\endcsname{0.9}
% scoring.d11_larva_band_min = 0.65  (mal-core/src/mal_core/scoring/scorers/d11_aquatic_structure.py:35)
\expandafter\def\csname cc@k@scoring.d11_larva_band_min\endcsname{0.65}
% scoring.d11_pupa_band_max = 0.15  (mal-core/src/mal_core/scoring/scorers/d11_aquatic_structure.py:36)
\expandafter\def\csname cc@k@scoring.d11_pupa_band_max\endcsname{0.15}
% scoring.d11_pupa_band_min = 0.02  (mal-core/src/mal_core/scoring/scorers/d11_aquatic_structure.py:36)
\expandafter\def\csname cc@k@scoring.d11_pupa_band_min\endcsname{0.02}
% scoring.d15_expansion_fraction = 0.2  (mal-core/src/mal_core/scoring/scorers/d15_persistence.py:125)
\expandafter\def\csname cc@k@scoring.d15_expansion_fraction\endcsname{0.2}
% scoring.d15_min_horizon_days = 30  (mal-core/src/mal_core/scoring/scorers/d15_persistence.py:97)
\expandafter\def\csname cc@k@scoring.d15_min_horizon_days\endcsname{30}
% scoring.d15_survival_fraction = 0.3  (mal-core/src/mal_core/scoring/scorers/d15_persistence.py:105)
\expandafter\def\csname cc@k@scoring.d15_survival_fraction\endcsname{0.3}
% scoring.d16_coverage_target = 0.5  (mal-core/src/mal_core/scoring/scorers/d16_detection_coverage.py:44)
\expandafter\def\csname cc@k@scoring.d16_coverage_target\endcsname{0.5}
% scoring.d16_radius_km_default = 5.0  (mal-core/src/mal_core/scoring/scorers/d16_detection_coverage.py:43)
\expandafter\def\csname cc@k@scoring.d16_radius_km_default\endcsname{5.0}
% scoring.d16_sensitivity_radius_far_km = 10.0  (mal-core/src/mal_core/scoring/scorers/d16_detection_coverage.py:45)
\expandafter\def\csname cc@k@scoring.d16_sensitivity_radius_far_km\endcsname{10.0}
% scoring.d16_sensitivity_radius_mid_km = 5.0  (mal-core/src/mal_core/scoring/scorers/d16_detection_coverage.py:45)
\expandafter\def\csname cc@k@scoring.d16_sensitivity_radius_mid_km\endcsname{5.0}
% scoring.d16_sensitivity_radius_near_km = 2.0  (mal-core/src/mal_core/scoring/scorers/d16_detection_coverage.py:45)
\expandafter\def\csname cc@k@scoring.d16_sensitivity_radius_near_km\endcsname{2.0}
% scoring.d19_activation_floor = 0.05  (mal-core/src/mal_core/scoring/scorers/d19_pool_persistence.py:81)
\expandafter\def\csname cc@k@scoring.d19_activation_floor\endcsname{0.05}
% scoring.d19_check_drying_weight = 0.2  (mal-core/src/mal_core/scoring/scorers/d19_pool_persistence.py:130)
\expandafter\def\csname cc@k@scoring.d19_check_drying_weight\endcsname{0.2}
% scoring.d19_check_floor_weight = 0.3  (mal-core/src/mal_core/scoring/scorers/d19_pool_persistence.py:134)
\expandafter\def\csname cc@k@scoring.d19_check_floor_weight\endcsname{0.3}
% scoring.d19_check_monotone_weight = 0.2  (mal-core/src/mal_core/scoring/scorers/d19_pool_persistence.py:132)
\expandafter\def\csname cc@k@scoring.d19_check_monotone_weight\endcsname{0.2}
% scoring.d19_check_persistence_weight = 0.3  (mal-core/src/mal_core/scoring/scorers/d19_pool_persistence.py:128)
\expandafter\def\csname cc@k@scoring.d19_check_persistence_weight\endcsname{0.3}
% scoring.d19_dry_days = 10  (mal-core/src/mal_core/scoring/scorers/d19_pool_persistence.py:84)
\expandafter\def\csname cc@k@scoring.d19_dry_days\endcsname{10}
% scoring.d19_pulse_days = 3  (mal-core/src/mal_core/scoring/scorers/d19_pool_persistence.py:78)
\expandafter\def\csname cc@k@scoring.d19_pulse_days\endcsname{3}
% scoring.d19_rain_pulse_mm = 30.0  (mal-core/src/mal_core/scoring/scorers/d19_pool_persistence.py:79)
\expandafter\def\csname cc@k@scoring.d19_rain_pulse_mm\endcsname{30.0}
% scoring.d1_band_max = 0.7  (mal-core/src/mal_core/scoring/scorers/d1_habitat_occupancy.py:33)
\expandafter\def\csname cc@k@scoring.d1_band_max\endcsname{0.7}
% scoring.d1_band_min = 0.2  (mal-core/src/mal_core/scoring/scorers/d1_habitat_occupancy.py:33)
\expandafter\def\csname cc@k@scoring.d1_band_min\endcsname{0.2}
% scoring.d1_center = 0.45  (mal-core/src/mal_core/scoring/scorers/d1_habitat_occupancy.py:34)
\expandafter\def\csname cc@k@scoring.d1_center\endcsname{0.45}
% scoring.d1_sigma = 0.15  (mal-core/src/mal_core/scoring/scorers/d1_habitat_occupancy.py:35)
\expandafter\def\csname cc@k@scoring.d1_sigma\endcsname{0.15}
% scoring.d20_cap_probe_mm = 80.0  (mal-core/src/mal_core/scoring/scorers/d20_washout_response.py:56)
\expandafter\def\csname cc@k@scoring.d20_cap_probe_mm\endcsname{80.0}
% scoring.d20_check_cap_weight = 0.2  (mal-core/src/mal_core/scoring/scorers/d20_washout_response.py:71)
\expandafter\def\csname cc@k@scoring.d20_check_cap_weight\endcsname{0.2}
% scoring.d20_check_no_washout_weight = 0.3  (mal-core/src/mal_core/scoring/scorers/d20_washout_response.py:69)
\expandafter\def\csname cc@k@scoring.d20_check_no_washout_weight\endcsname{0.3}
% scoring.d20_check_threshold_weight = 0.2  (mal-core/src/mal_core/scoring/scorers/d20_washout_response.py:73)
\expandafter\def\csname cc@k@scoring.d20_check_threshold_weight\endcsname{0.2}
% scoring.d20_check_washout_weight = 0.3  (mal-core/src/mal_core/scoring/scorers/d20_washout_response.py:67)
\expandafter\def\csname cc@k@scoring.d20_check_washout_weight\endcsname{0.3}
% scoring.d20_min_washout_fraction = 0.15  (mal-core/src/mal_core/scoring/scorers/d20_washout_response.py:48)
\expandafter\def\csname cc@k@scoring.d20_min_washout_fraction\endcsname{0.15}
% scoring.d20_no_washout_probe_mm = 30.0  (mal-core/src/mal_core/scoring/scorers/d20_washout_response.py:51)
\expandafter\def\csname cc@k@scoring.d20_no_washout_probe_mm\endcsname{30.0}
% scoring.d20_threshold_max_fraction = 0.01  (mal-core/src/mal_core/scoring/scorers/d20_washout_response.py:63)
\expandafter\def\csname cc@k@scoring.d20_threshold_max_fraction\endcsname{0.01}
% scoring.d20_threshold_probe_mm = 40.0  (mal-core/src/mal_core/scoring/scorers/d20_washout_response.py:61)
\expandafter\def\csname cc@k@scoring.d20_threshold_probe_mm\endcsname{40.0}
% scoring.d20_washout_probe_mm = 50.0  (mal-core/src/mal_core/scoring/scorers/d20_washout_response.py:46)
\expandafter\def\csname cc@k@scoring.d20_washout_probe_mm\endcsname{50.0}
% scoring.d21_horizon_day = 30  (mal-core/src/mal_core/scoring/scorers/d21_dispersal_kernel.py:112)
\expandafter\def\csname cc@k@scoring.d21_horizon_day\endcsname{30}
% scoring.d21_p50_band_max = 0.7  (mal-core/src/mal_core/scoring/scorers/d21_dispersal_kernel.py:39)
\expandafter\def\csname cc@k@scoring.d21_p50_band_max\endcsname{0.7}
% scoring.d21_p50_band_min = 0.15  (mal-core/src/mal_core/scoring/scorers/d21_dispersal_kernel.py:39)
\expandafter\def\csname cc@k@scoring.d21_p50_band_min\endcsname{0.15}
% scoring.d21_p50_center = 0.4  (mal-core/src/mal_core/scoring/scorers/d21_dispersal_kernel.py:41)
\expandafter\def\csname cc@k@scoring.d21_p50_center\endcsname{0.4}
% scoring.d21_p50_sigma = 0.2  (mal-core/src/mal_core/scoring/scorers/d21_dispersal_kernel.py:41)
\expandafter\def\csname cc@k@scoring.d21_p50_sigma\endcsname{0.2}
% scoring.d21_p90_band_max = 3.0  (mal-core/src/mal_core/scoring/scorers/d21_dispersal_kernel.py:40)
\expandafter\def\csname cc@k@scoring.d21_p90_band_max\endcsname{3.0}
% scoring.d21_p90_band_min = 0.7  (mal-core/src/mal_core/scoring/scorers/d21_dispersal_kernel.py:40)
\expandafter\def\csname cc@k@scoring.d21_p90_band_min\endcsname{0.7}
% scoring.d21_p90_center = 1.6  (mal-core/src/mal_core/scoring/scorers/d21_dispersal_kernel.py:42)
\expandafter\def\csname cc@k@scoring.d21_p90_center\endcsname{1.6}
% scoring.d21_p90_sigma = 0.8  (mal-core/src/mal_core/scoring/scorers/d21_dispersal_kernel.py:42)
\expandafter\def\csname cc@k@scoring.d21_p90_sigma\endcsname{0.8}
% scoring.d25_r_target = 0.6  (mal-core/src/mal_core/scoring/scorers/d25_cases_ghana.py:45)
\expandafter\def\csname cc@k@scoring.d25_r_target\endcsname{0.6}
% scoring.d2_lifetime_centre = 11.0  (mal-core/src/mal_core/scoring/scorers/d2_survival.py:70)
\expandafter\def\csname cc@k@scoring.d2_lifetime_centre\endcsname{11.0}
% scoring.d2_lifetime_width = 4.0  (mal-core/src/mal_core/scoring/scorers/d2_survival.py:70)
\expandafter\def\csname cc@k@scoring.d2_lifetime_width\endcsname{4.0}
% scoring.d2_window_max = 15  (mal-core/src/mal_core/scoring/scorers/d2_survival.py:47)
\expandafter\def\csname cc@k@scoring.d2_window_max\endcsname{15}
% scoring.d2_window_min = 8  (mal-core/src/mal_core/scoring/scorers/d2_survival.py:47)
\expandafter\def\csname cc@k@scoring.d2_window_min\endcsname{8}
% scoring.d3_band_max = 0.5  (mal-core/src/mal_core/scoring/scorers/d3_eip.py:69)
\expandafter\def\csname cc@k@scoring.d3_band_max\endcsname{0.5}
% scoring.d3_band_min = 0.2  (mal-core/src/mal_core/scoring/scorers/d3_eip.py:69)
\expandafter\def\csname cc@k@scoring.d3_band_min\endcsname{0.2}
% scoring.d3_center = 0.35  (mal-core/src/mal_core/scoring/scorers/d3_eip.py:72)
\expandafter\def\csname cc@k@scoring.d3_center\endcsname{0.35}
% scoring.d3_day = 30  (mal-core/src/mal_core/scoring/scorers/d3_eip.py:67)
\expandafter\def\csname cc@k@scoring.d3_day\endcsname{30}
% scoring.d3_sigma = 0.15  (mal-core/src/mal_core/scoring/scorers/d3_eip.py:72)
\expandafter\def\csname cc@k@scoring.d3_sigma\endcsname{0.15}
% scoring.d4_collapse_max = 0.8  (mal-core/src/mal_core/scoring/scorers/d4_stability.py:70)
\expandafter\def\csname cc@k@scoring.d4_collapse_max\endcsname{0.8}
% scoring.d4_collapse_min = 0.1  (mal-core/src/mal_core/scoring/scorers/d4_stability.py:70)
\expandafter\def\csname cc@k@scoring.d4_collapse_min\endcsname{0.1}
% scoring.d4_explosion_max = 5.0  (mal-core/src/mal_core/scoring/scorers/d4_stability.py:77)
\expandafter\def\csname cc@k@scoring.d4_explosion_max\endcsname{5.0}
% scoring.d5_band_max = 0.85  (mal-core/src/mal_core/scoring/scorers/d5_morans.py:95)
\expandafter\def\csname cc@k@scoring.d5_band_max\endcsname{0.85}
% scoring.d5_band_min = 0.2  (mal-core/src/mal_core/scoring/scorers/d5_morans.py:95)
\expandafter\def\csname cc@k@scoring.d5_band_min\endcsname{0.2}
% scoring.d5_center = 0.5  (mal-core/src/mal_core/scoring/scorers/d5_morans.py:98)
\expandafter\def\csname cc@k@scoring.d5_center\endcsname{0.5}
% scoring.d5_sigma = 0.25  (mal-core/src/mal_core/scoring/scorers/d5_morans.py:98)
\expandafter\def\csname cc@k@scoring.d5_sigma\endcsname{0.25}
% scoring.d8_min_score = 0.4  (mal-core/src/mal_core/scoring/scorers/d8_coupling.py:35)
\expandafter\def\csname cc@k@scoring.d8_min_score\endcsname{0.4}
% scoring.d8_pass_threshold = 0.3  (mal-core/src/mal_core/scoring/scorers/d8_coupling.py:76)
\expandafter\def\csname cc@k@scoring.d8_pass_threshold\endcsname{0.3}
% scoring.d8_score_divisor = 0.5  (mal-core/src/mal_core/scoring/scorers/d8_coupling.py:70)
\expandafter\def\csname cc@k@scoring.d8_score_divisor\endcsname{0.5}
% scoring.gate_census_reference = 27000000  (mal-core/src/mal_core/scoring/scorers/gates.py:191)
\expandafter\def\csname cc@k@scoring.gate_census_reference\endcsname{27000000}
% scoring.gate_ratio_band_max = 1.2  (mal-core/src/mal_core/scoring/scorers/gates.py:214)
\expandafter\def\csname cc@k@scoring.gate_ratio_band_max\endcsname{1.2}
% scoring.gate_ratio_band_min = 0.8  (mal-core/src/mal_core/scoring/scorers/gates.py:214)
\expandafter\def\csname cc@k@scoring.gate_ratio_band_min\endcsname{0.8}
% scoring.gate_row_sum_threshold = 0.99  (mal-core/src/mal_core/scoring/scorers/gates.py:271)
\expandafter\def\csname cc@k@scoring.gate_row_sum_threshold\endcsname{0.99}
% scoring.gate_row_sum_tolerance = 0.01  (mal-core/src/mal_core/scoring/scorers/gates.py:260)
\expandafter\def\csname cc@k@scoring.gate_row_sum_tolerance\endcsname{0.01}
% scoring.gate_rural_class = 50  (mal-core/src/mal_core/scoring/scorers/gates.py:317)
\expandafter\def\csname cc@k@scoring.gate_rural_class\endcsname{50}
% scoring.gate_sample_files = 48  (mal-core/src/mal_core/scoring/scorers/gates.py:59)
\expandafter\def\csname cc@k@scoring.gate_sample_files\endcsname{48}
% scoring.gate_urban_class = 30  (mal-core/src/mal_core/scoring/scorers/gates.py:316)
\expandafter\def\csname cc@k@scoring.gate_urban_class\endcsname{30}
% scoring.gate_urban_ratio_max = 0.7  (mal-core/src/mal_core/scoring/scorers/gates.py:336)
\expandafter\def\csname cc@k@scoring.gate_urban_ratio_max\endcsname{0.7}
% scoring.gate_wall_time_limit_s = 30.0  (mal-core/src/mal_core/scoring/scorers/gates.py:158)
\expandafter\def\csname cc@k@scoring.gate_wall_time_limit_s\endcsname{30.0}
% scoring.regression_threshold = -0.05  (mal-core/src/mal_core/scoring/feedback.py:51)
\expandafter\def\csname cc@k@scoring.regression_threshold\endcsname{-0.05}
% scoring.verdict_borderline = 0.5  (mal-core/src/mal_core/scoring/runner.py:24)
\expandafter\def\csname cc@k@scoring.verdict_borderline\endcsname{0.5}
% scoring.verdict_regressed = 0.3  (mal-core/src/mal_core/scoring/runner.py:24)
\expandafter\def\csname cc@k@scoring.verdict_regressed\endcsname{0.3}
% scoring.verdict_viable = 0.7  (mal-core/src/mal_core/scoring/runner.py:24)
\expandafter\def\csname cc@k@scoring.verdict_viable\endcsname{0.7}
% training.ckpt_period = 10  (mal-core/src/mal_core/training/trainer.py:102)
\expandafter\def\csname cc@k@training.ckpt_period\endcsname{10}
% training.dice_count_threshold = 0.5  (mal-core/src/mal_core/training/model.py:110)
\expandafter\def\csname cc@k@training.dice_count_threshold\endcsname{0.5}
% training.dice_eps = 1e-06  (mal-core/src/mal_core/training/model.py:97)
\expandafter\def\csname cc@k@training.dice_eps\endcsname{1e-06}
% training.dice_tau = 0.5  (mal-core/src/mal_core/training/model.py:87)
\expandafter\def\csname cc@k@training.dice_tau\endcsname{0.5}
% training.dice_thresh = 0.5  (mal-core/src/mal_core/training/model.py:87)
\expandafter\def\csname cc@k@training.dice_thresh\endcsname{0.5}
% training.dice_weight = 0.5  (mal-core/src/mal_core/training/model.py:105)
\expandafter\def\csname cc@k@training.dice_weight\endcsname{0.5}
% training.env_channels = 4  (mal-core/src/mal_core/training/dataset.py:22)
\expandafter\def\csname cc@k@training.env_channels\endcsname{4}
% training.fallback_in_channels = 6  (mal-core/src/mal_core/training/wrapper.py:62)
\expandafter\def\csname cc@k@training.fallback_in_channels\endcsname{6}
% training.fallback_out_channels = 2  (mal-core/src/mal_core/training/wrapper.py:66)
\expandafter\def\csname cc@k@training.fallback_out_channels\endcsname{2}
% training.pad_multiple = 16  (mal-core/src/mal_core/training/wrapper.py:100)
\expandafter\def\csname cc@k@training.pad_multiple\endcsname{16}
% training.patch_size = 128  (mal-core/src/mal_core/training/dataset.py:19)
\expandafter\def\csname cc@k@training.patch_size\endcsname{128}
% training.sub_seed = 42  (mal-core/src/mal_core/training/dataset.py:210)
\expandafter\def\csname cc@k@training.sub_seed\endcsname{42}
% training.transmission_channels = 4  (mal-core/src/mal_core/training/dataset.py:21)
\expandafter\def\csname cc@k@training.transmission_channels\endcsname{4}
% training.unet_channels = 32, 64, 128, 256  (mal-core/src/mal_core/training/model.py:46)
\expandafter\def\csname cc@k@training.unet_channels\endcsname{32, 64, 128, 256}
% training.vector_channels = 2  (mal-core/src/mal_core/training/dataset.py:20)
\expandafter\def\csname cc@k@training.vector_channels\endcsname{2}


\begin{document} \newcommand{\src}[1]{\code{#1}}

\title{The \emph{ingest} stage:\\ turning downloaded layers into the environment the engine can
read} \author{MalariaSentinel conceptual documentation} \date{}

\begin{abstract} \noindent The \emph{ingest} stage is the transform layer of the MalariaSentinel
pipeline. It takes raw, downloaded, provider-native layers --- rainfall, air temperature,
permanent-surface water, land cover, a coastline, and a family of host-density layers --- and
turns them into the three artefacts the C++ simulation engine consumes: a daily environment
NetCDF, a static habitat-patch GeoPackage and a static host grid. It
does not fetch data on the environment path, it does not simulate, and it does not score. This
document explains what the stage computes, in what order, on which grid, and what its numbers
mean and do not mean. Every constant quoted here is generated from the shared constants file or
from the generated tables; every behavioural claim carries a file-and-line citation.
\end{abstract}

\tableofcontents

\bigskip \noindent\textbf{Notation.} Symbols used in this document are named on first use, with
the code symbol that implements them in parentheses, e.g.\ the static habitat potential $H[i,g]$
(\code{static\_potential} in \src{mal-core/src/mal\_core/habitat/potential.py:50}). File
citations are given as \src{path:line} or \src{path:line--line} and can be checked by opening
the file at that line.

\section{Purpose, and the boundary of the purpose} \label{sec:purpose}

\subsection{What question the stage answers}

The engine that simulates malaria transmission needs, for every cell of a fixed grid and for
every day of the simulated period, a small set of numbers it can read without interpreting
satellite imagery or provider formats:

\begin{itemize}
\item how much rain fell, and how warm it was, on each day; \item how much standing-water
breeding habitat each cell statically offers, expressed as an index in $[0,1]$ per habitat type;
\item which cells are candidate breeding patches at all; \item where the people are, where the
livestock are, and how each of those two populations moves between cells.
\end{itemize}

The \emph{ingest} stage answers exactly that question, on the AOI's own grid, and nothing wider.
It is a transform with a fixed output contract: given an area of interest and the raw products
registered for it, it emits a self-describing set of files on disk and registers them in the
central manifest. The two public builders are \code{build\_env\_tensor} (the daily environment
NetCDF and the habitat GeoPackage) and \code{build\_host\_dataset} (the static host grid)
\src{mal-core/src/mal\_core/ingest/\_\_init\_\_.py}.

\subsection{What would count as the stage having answered it}

The stage has answered its question when, for one AOI, the following are all true:

\begin{enumerate}
\item the environment NetCDF exists, carries the daily channels and the static planes on one
grid, and declares its contract version (\S\ref{sec:envnc}); 
\item the habitat GeoPackage exists
and enumerates the candidate cells with their per-type potential and their grid identity
(\S\ref{sec:gpkg}); 
\item the host grid exists with one plane per host layer, all on the same
grid (\S\ref{sec:hosts}); 
\item every one of those artefacts is registered in the AOI manifest
under its own key (\S\ref{sec:manifest}).
\end{enumerate}

\subsection{What the stage is explicitly \emph{not} for}

This list is a filter, not a courtesy. If a reader is looking for the following, it is
elsewhere.

\begin{itemize}
\item \textbf{It is not the download stage.} The environment path fetches nothing: it resolves
already-downloaded files through the manifest (\S\ref{sec:envinputs}). The host
path \emph{does} fetch, through the provider loaders it calls directly
(\S\ref{sec:hostdownloads}) --- an asymmetry the stage does not hide. 
\item \textbf{It is not
the simulator.} No mosquito is drawn and no day is stepped. The dynamic parts of the habitat
model --- the two temporal responses, the daily accuracy, the daily capacity --- are computed by
the engine, every day, from the static planes this stage writes (\S\ref{sec:dynamic}). 
\item
\textbf{It is not the scorer, the trainer or the predictor.} The generated contract table
(\S\ref{sec:contracts}) lists \code{scoring}, \code{prediction} and \code{abm} as consumers of
the environment NetCDF. 
\item \textbf{It is not a measurement device.} The habitat numbers this
stage produces are indices of relative offer: not pool counts, not water depths, not volumes,
not probabilities (\S\ref{sec:hnot}). 
\item \textbf{It is not a terrain analysis.} No digital
elevation model, no topographic wetness index, no catchment computation, no millimetre water
balance (\S\ref{sec:absences}). 
\item \textbf{It does not choose a species}; which affinity row
applies is decided downstream.
\end{itemize}

\subsection{Which authority governs which statement} \label{sec:authority}

This document owns the description of the \emph{ingest} stage, and it settles its own authority
questions; there is no shared page that arbitrates on its behalf. Three authorities apply, and
each governs a different kind of claim. \emph{Values} --- the numbers the pipeline uses --- are
governed by the shared constants file \code{mal-core/config/model\_constants.yaml}, and are
quoted here through generated macros that read it; a value with no constants-file key, because it lives in the code
instead, is quoted through the generated code-constant macros, which read the implementation
itself. \emph{Contract values and names} --- the contract version, the manifest keys, the
artefact names --- are governed by the generated contract and parameter tables
(\S\ref{sec:contracts}), which are scanned from the code rather than restated by hand.
\emph{Semantics} --- what a number means, in what order the steps run, what a step deliberately
does not do --- are governed by this document, because this document is the one that owns the
ingest content. Where two of these disagree, the lower authority is the bug: the constants file
and the code are not adjusted to fit a sentence, and the sentence is corrected. That is also why
the numbers this document quotes are generated macros rather than literals, and why a value the
code changes makes this document's build fail instead of making the document lie.

\section{Entry points, artefacts, and the manifest keys they claim} \label{sec:entry}

\subsection{Entry points}

The stage is reached in two ways.

\begin{description}
\item[The command line.] \code{malariasim ingest} is a subcommand with the options \code{--aoi},
\code{--year}, \code{--month}, \code{--output-dir}, \code{--data-dir} and \code{--what}
\src{mal-core/src/mal\_core/cli.py:198--206}. The switch \code{--what} selects which of the
two builders runs; \code{all} is the default \src{mal-core/src/mal\_core/cli.py:196--239}.
\item[Python.] The two builders are importable and independent
\src{mal-core/src/mal\_core/ingest/\_\_init\_\_.py}.
\end{description}

\subsection{Artefacts and their manifest keys}

The identity of every artefact this stage produces is its manifest key. Those keys live in
exactly one table, the ingest-artefact table (\code{INGEST\_ARTIFACTS}), which is the only new
table in the dataset catalogue \src{mal-core/src/mal\_core/download/catalog.py:55--92}.

\begin{center}
\small
\begin{longtable}{@{}llllr@{}}
\toprule
Manifest key & Producer & On-disk name & Format & Required for the engine \\
\midrule
\endhead
\code{env} & \code{build\_env\_tensor} & \code{<aoi>\_regional\_<y0>\_<y1>\_env.nc} & NetCDF &
yes \\ \code{habitat} & \code{build\_env\_tensor} & \code{<aoi>\_habitat\_patches.gpkg} &
GeoPackage & yes \\ \code{host\_static} & \code{build\_host\_dataset} &
\code{<aoi>\_host\_static.nc} & NetCDF & yes \\ \code{host\_manifest} &
\code{build\_host\_dataset} & \code{host\_manifest.json} & JSON & no \\
\bottomrule
\end{longtable}
\end{center}

\smallskip \noindent The \emph{required for the engine} column is the catalogue's own
declaration \src{mal-core/src/mal\_core/download/catalog.py:55--92}. It is worth noting that the
\emph{manifest entry} written for the environment NetCDF and the habitat GeoPackage does
\emph{not} carry that flag: those two registrations do not pass \code{required\_for\_abm}
\src{mal-core/src/mal\_core/ingest/env.py:147--162},
\src{mal-core/src/mal\_core/ingest/env.py:174--187}, and the manifest writer therefore stores
the default \code{False} \src{mal-core/src/mal\_core/download/manifest.py:118--189}. The
verified manifest for Ghana shows \code{required\_for\_abm: false} on both entries while the
host grid carries \code{true}. The completeness check the engine
performs before it runs (\S\ref{sec:seam}) therefore does not police the environment artefacts
on this flag.

\subsection{Where the CLI puts things, and where the manifest says they are}

The CLI's defaults for the two builders are not mutually consistent, and the inconsistency is
observable rather than theoretical. The environment builder is called with the output directory
serving as \emph{both} input directory and output directory, and with the data root taken as the
parent of that directory \src{mal-core/src/mal\_core/cli.py:221--222}. The host builder is
called with the output directory only \src{mal-core/src/mal\_core/cli.py:225}. Meanwhile the host builder registers its artefact with
the default data root, not with the directory it wrote to
\src{mal-core/src/mal\_core/ingest/hosts.py:219--240}.

The consequence, stated as a described situation: an operator runs the documented command
\code{malariasim ingest --aoi ghana --what env} with the default output directory. The
environment builder looks for the manifest in its own output directory rather than in the AOI's
data directory, and the build fails immediately with a missing-required-input error naming the
download command that would fix it \src{mal-core/src/mal\_core/ingest/daily\_nc.py:186--190},
\src{mal-core/src/mal\_core/download/catalog.py:176--268}. A run that does succeed can leave the
NetCDF in one directory, the host grid in another, and the manifest entries naming files that
are not where the manifest says they are --- which the completeness validator reports as missing
\src{mal-core/src/mal\_core/download/manifest.py:192--245}.

\section{The spatial frame} \label{sec:frame}

\subsection{The area of interest is one of the frame's two owners}

The artefacts this stage writes live on one grid, and that grid has two owners rather than one:
the environment NetCDF's shape is inherited from the rainfall file the builder reads first, and
only the habitat GeoPackage's shape comes from the area-of-interest record (both computations
are set out below). The record is a bounding box, a coordinate reference system, a nominal
ground resolution in metres, and a scale label
\src{mal-commonlib/src/mal\_commonlib/aoi.py:64--95}. The registry that supplies those values
for each slug is a single YAML file \src{mal-commonlib/src/mal\_commonlib/aois.yaml:16--23}; for
the Ghana slug it declares a bounding box of
$[\codeconst{ingest.bbox_west}, \codeconst{ingest.bbox_south},
\codeconst{ingest.bbox_east}, \codeconst{ingest.bbox_north}]$, a geographic CRS, a
nominal resolution of \codeconst{ingest.cell_size_m} metres, and a regional scale. The
bounding box and the resolution are declared literally in the registry
\src{mal-commonlib/src/mal\_commonlib/aois.yaml:18--20}, and are quoted here through the
generated code-constant macros, which read the registry itself. The resolution
\codeconst{ingest.cell_size_m} is the AOI record's own nominal ground resolution; it is a
different quantity from \mskey{ingest.cell_size_hint_m}, which is the fallback cell size used
when the rainfall NetCDF carries no \code{resolution\_m} attribute
\src{mal-core/src/mal\_core/ingest/daily\_nc.py:243--247}. The two have the same magnitude
today, but they are independent literals: the AOI record never reads the constants file.

The grid size is derived, not declared: the metric extent of the bounding box is divided by the
nominal resolution and rounded up to whole cells
\src{mal-commonlib/src/mal\_commonlib/aoi.py:195--199}. The metric extent is itself an
approximation --- degrees are converted to metres at the bounding-box centroid latitude
\src{mal-commonlib/src/mal\_commonlib/aoi.py:154--171}. So the frame is defined in metres but
\emph{written} in the CRS's own units, and for a geographic CRS those units are degrees.

\subsection{Two owners of the grid shape}

There are two computations of the grid height and width in this stage, and they must agree.

\begin{enumerate}
\item The environment NetCDF takes its shape from the file it reads first, the rainfall NetCDF:
\code{n\_days, h, w = rainfall\_raw.shape} \src{mal-core/src/mal\_core/ingest/daily\_nc.py:203}.
The rainfall file was written by the download stage, already reduced to the AOI grid. 
\item The
habitat GeoPackage takes its shape from the AOI record: \code{h, w = aoi.cells\_per\_side()}
\src{mal-core/src/mal\_core/ingest/env.py:69}, and the static planes are re-read from the NetCDF
with a shape assertion against that same \code{cells\_per\_side()}
\src{mal-core/src/mal\_core/ingest/env.py:42--44}.
\end{enumerate}

The link between the two is that assertion: if the NetCDF's statics do not have the shape the
AOI record implies, the GeoPackage build raises
\src{mal-core/src/mal\_core/ingest/env.py:43--44}. Nothing, however, checks that the rainfall
file the environment builder read was itself built on \emph{this} AOI's grid; the frame is thus
owned jointly by the download stage and the AOI record, with only a shape check between them
(\S\ref{sec:open} question 13).

\subsection{What the coordinates in the environment NetCDF actually are}

The NetCDF inherits the coordinate vectors of the rainfall file
\src{mal-core/src/mal\_core/ingest/daily\_nc.py:200--201}. Whether those vectors carry real
coordinates or merely cell indices depends on how the source file was written, and the stage
does not decide.

The builder adds CF-1.8 two-dimensional latitude and longitude variables and a grid-mapping
variable only under a condition: that the median absolute spacing of both axes is a positive
number of degrees below \mskey{ingest.grid_degrees_step_max}
\src{mal-core/src/mal\_core/ingest/daily\_nc.py:467--494}. Verified on the real Ghana inputs,
the rainfall file declares dimensions \code{y} and \code{x} with no coordinate variables, so the
inherited axes are integer cell indices; the median spacing is therefore one cell, which does
not satisfy the condition, and the latitude, longitude and \code{crs} variables are not written.
The committed NetCDF \code{data/ghana/ghana\_regional\_2024\_2024\_env.nc} confirms this: its
variables are the \codeconst{ingest.env_daily_channels} daily channels, the
\codeconst{ingest.env_static_planes} static planes, and the \code{time}, \code{y} and
\code{x} coordinates, with no \code{lat}, \code{lon} or \code{crs} variable and no coordinate
reference system of any kind. In other words, in the delivered artefact the environment NetCDF
is a bare array; the geolocation is supplied to the engine separately, as the bounding box,
resolution and scale of the AOI forwarded on the command line
\src{mal-core/src/mal\_core/abm/wrapper.py:199--202}.

\subsection{Cell area is an approximation, and the magnitude is stated}

The engine needs the ground area of a cell to turn a dimensionless accuracy into an absolute
capacity. Both the Python reference module and the C++ engine compute it with the same
equirectangular approximation: the north--south side is taken to be exactly the nominal
resolution, and the east--west side is that resolution times the cosine of the latitude
\src{mal-core/src/mal\_core/habitat/capacity.py:39--46},
\src{mal-core/src/mal\_core/abm/include/mal\_abm\_fast/habitat\_availability.hpp:188--194}. The
engine applies it per grid row, at the row's centre latitude
\src{mal-core/src/mal\_core/abm/src/coordinator.cpp:130--134}.

This is an approximation and it is biased. The true area of a cell of the grid actually written
--- a bounding box in degrees, cut into a whole number of cells --- differs from the
approximation for two reasons: the cell count is rounded up, so the cell edges do not fall on
whole kilometres; and the north--south extent of the written grid is a constant number of
degrees rather than a constant number of metres. The code computes the approximation only: the
size of the difference for the Ghana box was measured outside the codebase during the
preparation of this document --- roughly eight tenths of a percent, with the approximation the
smaller of the two --- and it is listed under ``Not verified, and not claimed''
(\S\ref{sec:appendix}); it is not a fact this codebase asserts. Because both the Python module
and the C++ engine use the same formula, the two mirrors agree with each other exactly; the
discrepancy is a small systematic bias in every capacity number, proportional to cell area, not
a disagreement between implementations.

The Python module states the factor explicitly: edge cells clipped by the AOI bounding box are
not area-corrected, and per-row accuracy is a documented sensitivity extension
\src{mal-core/src/mal\_core/habitat/capacity.py:1--14}.

\section{Inputs: what is required, what is optional, and how it is resolved}
\label{sec:envinputs}

\subsection{The environment path is manifest-only}

The environment builder declares its entire dataset knowledge as a table of slots, each
identified by a manifest key \src{mal-core/src/mal\_core/ingest/daily\_nc.py:87--98}. Resolution
goes through the catalogue, which reads the AOI's manifest (\code{manifest.json}) and treats it
as the single source of truth \src{mal-core/src/mal\_core/download/catalog.py:176--268}. There
is no conventional-filename fallback inside the consumer: a dataset that was never registered
does not exist as far as this stage is concerned
\src{mal-core/src/mal\_core/download/manifest.py:58--68}.

\begin{center}
\small
\begin{longtable}{@{}lll@{}}
\toprule
Slot & Manifest key & Requirement \\
\midrule
\endhead
\code{rainfall} & \code{chirps\_rainfall\_daily} & required \\ \code{air\_temp} &
\code{era5\_air\_temp\_daily} & required \\ \code{jrc\_features} & \code{jrc\_habitat\_features}
& required \\ \code{wc\_fractions} & \code{dw\_class\_fractions} & required \\ \code{land\_mask}
& \code{coastline\_land\_mask} & optional \\ \code{host\_static} & \code{host\_static} &
optional \\ \code{salinity} & \code{smap\_salinity} & optional \\
\bottomrule
\end{longtable}
\end{center}

\smallskip \noindent A missing required slot is a hard error that names the download command
that would produce it \src{mal-core/src/mal\_core/download/catalog.py:255--263}. A manifest
entry that exists but whose file is missing on disk is also a hard error --- a stale manifest is
never tolerated silently \src{mal-core/src/mal\_core/download/manifest.py:108--114}. A missing
optional slot yields \code{None}, and the corresponding behaviour is skipped.

\subsection{The declared \code{host\_static} slot is dead}

The environment builder declares an optional dependency on the host grid
\src{mal-core/src/mal\_core/ingest/daily\_nc.py:96}, and then never reads it: the resolved value
is not referenced anywhere in the builder body
\src{mal-core/src/mal\_core/ingest/daily\_nc.py:186--260}. The declaration is
a dead dependency, recorded here rather than hidden (\S\ref{sec:open} question 12).

\subsection{What each input is, natively, and what the reduction costs}

The following table is the stage's input inventory. The native resolution at which each
provider publishes its product is a property of the download stage's loaders and is tabulated
in the generated table of the loaders' native grids (\S\ref{sec:contracts}); it is not restated
here. The target resolution is always the AOI grid of \S\ref{sec:frame}.

\begin{center}
\small
\begin{center}
\small
\begin{longtable}{@{}p{3.0cm}p{2.8cm}p{1.9cm}p{6.9cm}@{}}
\toprule
Input & Provider & Target & Units; what the reduction loses \\
\midrule
\endhead
\code{chirps\_rainfall\_daily} & CHIRPS v2.0 & AOI grid
  & mm/day. Nearest-neighbour reprojection
\src{mal-commonlib/src/mal\_commonlib/data/loaders/chirps.py:230--242}. Loses the sub-cell
structure of each day. \\ \code{era5\_air\_temp\_daily} & ERA5-Land & AOI grid
  & $\!{}^\circ$C daily mean 2\,m air temperature
\src{mal-commonlib/src/mal\_commonlib/data/loaders/era5.py:107}. Loses the sub-cell gradient. \\
\code{jrc\_habitat\_features} & JRC GSW v1.4 & AOI grid
  & Fraction in $[0,1]$ of the cell's pixels with water occurrence at least
\codeconst{ingest.jrc_occurrence_threshold_pct}\,\%; the fraction, not the retired binary mask
\src{mal-commonlib/src/mal\_commonlib/data/loaders/jrc\_gsw.py:7--12}. Loses the position of the
water inside the cell. \\ \code{dw\_class\_fractions} & Dynamic World V1 & AOI grid
  & Four areal fractions in $[0,1]$: built, flooded vegetation, crops, and
open-plus-water \src{mal-commonlib/src/mal\_commonlib/data/loaders/dynamic\_world.py:94--103};
per year a modal label reduced to an areal fraction
\src{mal-commonlib/src/mal\_commonlib/data/loaders/dynamic\_world.py:273--334}. Loses the
arrangement of classes inside the cell. \\ \code{coastline\_land\_mask} & GSHHG full & AOI grid
  & Land fraction in $[0,1]$ per cell
\src{mal-commonlib/src/mal\_commonlib/data/loaders/coastline.py:1--40}. Loses the shape of the
shoreline inside the cell. \\ \code{smap\_salinity} & SMAP RSS L3 (monthly) & AOI grid, bilinear
  & Practical salinity units; reprojected onto the environment grid by the
ingest stage itself \src{mal-core/src/mal\_core/ingest/daily\_nc.py:550--578}. \\
\bottomrule
\end{longtable}
\end{center}
\end{center}

\smallskip \noindent Two inputs are resampled by the ingest stage itself. The coastline
GeoTIFF is read and resampled bilinearly when its shape differs from the environment grid
\src{mal-core/src/mal\_core/ingest/daily\_nc.py:126--148}, and the monthly salinity field is
reprojected onto the environment grid, one month at a time, with a bilinear resampling
\src{mal-core/src/mal\_core/ingest/daily\_nc.py:550--578}. Everything else arrives already on
the AOI grid.

\subsection{The host path downloads} \label{sec:hostdownloads}

The host builder does not go through the manifest. It calls the provider loaders directly:
WorldPop population \src{mal-core/src/mal\_core/ingest/hosts.py:74}, GLW4 livestock
\src{mal-core/src/mal\_core/ingest/hosts.py:92}, GHS-SMOD urban class
\src{mal-core/src/mal\_core/ingest/hosts.py:111}, Overture building footprints
\src{mal-core/src/mal\_core/ingest/hosts.py:146}, and the derived wildlife host proxy
\src{mal-core/src/mal\_core/ingest/hosts.py:169}. Each of those loaders fetches from its
provider when its cache is cold. So a ``run ingest'' on a machine with an empty cache performs
network transfers on the host path, and performs none on the environment path. The retired
specification's claim that ingest never downloads holds only for the environment path
(\S\ref{sec:open} question 1).

\section{The environment build: order of operations and its consequences} \label{sec:order}

\subsection{The order, as pseudocode}

\begin{verbatim}
build_daily_env_nc(aoi, data_dir, output_dir):
  resolve the four required slots and the three optional slots
    through the manifest; a missing required slot is fatal

  rain  <- open(rainfall file)          # (time, y, x); gives h, w and the time axis
  fill the rainfall NoData per day from the nearest always-valid cell
  clamp negative rainfall to zero

  if a coastline file is available:
      land_frac   <- read it, resampled to (h, w)
      dist        <- distance to land, in metres
      ocean_mask  <- (land_frac < sea_water_floor) AND (dist > coast_band_m)
      set rainfall to zero on the ocean mask
  else:
      print a warning; the saltwater filter is skipped

  temp <- open(air temperature file)
  raise unless the temperature series covers every day of the rain axis
  re-index the temperature onto the rain axis
  set temperature to zero on the ocean mask

  read the static planes; average a multi-year stack over its year axis
  zero all static planes on the ocean mask
  print the coastline census

  H        <- static_potential(five planes)
  H_total  <- total_potential(H)
  census the cells whose H_total exceeds the ingest floor; raise if empty

  if a salinity file is available:
      reproject each month onto the grid; broadcast each month onto its days

  assemble the dataset: the daily channels, the static planes, optional land_frac
    and salinity_ppt; set the global attributes
  add CF latitude/longitude/crs variables if the axes are a degree grid
  write the NetCDF

  (in build_env_tensor) re-read the static planes out of the NetCDF that was
    just written, derive the habitat GeoPackage from them, register both
    artefacts in the manifest
\end{verbatim}

\smallskip \noindent The dataset assembled by the last step is the one the NetCDF contract
states: \codeconst{ingest.env_daily_channels} daily channels and
\codeconst{ingest.env_static_planes} static planes, plus the optional \code{land\_frac} and
\code{salinity\_ppt} variables \src{mal-core/src/mal\_core/ingest/daily\_nc.py:404--431}.

\subsection{The order as a table}

\begin{center}
\small
\begin{longtable}{@{}p{0.6cm}p{4.4cm}p{3.4cm}p{5.4cm}@{}}
\toprule
\# & Step & Actor & Why here \\
\midrule
\endhead
1 & resolve slots & ingest & every later step needs a path; the manifest is the only identity \\
2 & read rainfall & ingest & fixes the time axis, the day count and the grid shape \\ 3 & fill
rainfall NoData & ingest & must precede the ocean mask, which reads the same array \\ 4 &
coastline and open-sea mask & ingest & needed before the statics are zeroed and before
temperature is masked \\ 5 & read and re-index temperature & ingest & the rain axis is already
fixed, so the coverage check is possible \\ 6 & mask temperature on the ocean & ingest & follows
the mask's construction \\ 7 & assemble the static planes & ingest & the multi-year average must
precede the mask application \\ 8 & zero the statics on the ocean & ingest & the mask exists;
the statics are the last thing touched \\ 9 & censuses & ingest & an auditable record of what
the mask did \\ 10 & compute $H$ and the candidate census & ingest & needs the final statics \\
11 & salinity broadcast & ingest & after the mask, so the broadcast is a value on the final grid
\\ 12 & assemble, georeference, write & ingest & one write; the file is the contract \\ 13 &
derive the GeoPackage & ingest & from the file just written, so the two cannot disagree \\
\bottomrule
\end{longtable}
\end{center}

\subsection{Consequences worth stating}

\begin{enumerate}
\item \textbf{The rainfall axis owns the time axis.} The temperature series is re-indexed onto
it, so the environment NetCDF's days are exactly the rainfall file's days, and a temperature
file that does not cover them is a loud failure rather than a silent broadcast
\src{mal-core/src/mal\_core/ingest/daily\_nc.py:262--284}. The retired path, in which an annual
temperature raster was broadcast across every day, no longer exists
\src{mal-core/src/mal\_core/ingest/daily\_nc.py:257--261}. 
\item \textbf{The absence of the
coastline artefact degrades the build instead of failing it.} The saltwater filter is skipped
with a printed note, and the static planes keep their open-ocean cells
\src{mal-core/src/mal\_core/ingest/daily\_nc.py:357--362}. The stated purpose of that filter is
precisely to stop open-ocean cells from becoming breeding candidates
\src{mal-core/src/mal\_core/ingest/daily\_nc.py:327--339} (\S\ref{sec:open} question 11). 
\item
\textbf{The GeoPackage cannot disagree with the NetCDF's statics.} The GeoPackage is derived by
re-reading the static planes out of the NetCDF that was written first
\src{mal-core/src/mal\_core/ingest/env.py:143--173}. The consequence runs the other way too: the
GeoPackage's candidate set is computed from the already-filtered statics, so it inherits the
ocean mask. 
\item \textbf{Order of the two censuses.} The coastline census is printed before the
candidate census \src{mal-core/src/mal\_core/ingest/daily\_nc.py:364--387}, so a reader of the
log sees how many cells the mask removed before seeing how many candidates survived.
\end{enumerate}

\section{The two daily channels} \label{sec:channels}

\subsection{Rainfall: units, donor fill, and the assumption behind it}

The rainfall channel is CHIRPS v2.0 daily precipitation, in millimetres per day
\src{mal-core/src/mal\_core/ingest/daily\_nc.py:406--408}. Its NoData handling is a spatial
donor fill, applied per day.

A cell is \emph{eligible} to donate when its fraction of NoData days over the whole series is
below \mskey{ingest.fill_frac_valid_max}
\src{mal-core/src/mal\_core/ingest/daily\_nc.py:209--212}. For every day, each NoData pixel is
replaced by the value of the spatially nearest eligible cell, found by a Euclidean distance
transform over the eligible mask \src{mal-core/src/mal\_core/ingest/daily\_nc.py:213--220}. If
no cell is eligible at all, the NoData becomes zero
\src{mal-core/src/mal\_core/ingest/daily\_nc.py:221--222}. Negative precipitation, which some
gauge products report, is clamped to zero \src{mal-core/src/mal\_core/ingest/daily\_nc.py:223}.

The assumption the code makes and does not test is that the nearest eligible cell is a plausible
donor for that day. The choice is a \emph{choice}: a temporal interpolation would have been the
alternative. The stage does not record which cells were filled beyond a count in the build
summary \src{mal-core/src/mal\_core/ingest/daily\_nc.py:504--514}.

\subsection{Air temperature: units, coverage, and the zero sentinel}

The temperature channel is the ERA5-Land daily mean 2\,m air temperature, in degrees Celsius
\src{mal-core/src/mal\_core/ingest/daily\_nc.py:409--411}. The name of the channel says what it
measures: the retired annual \code{water\_temp\_c} stack is not read
\src{mal-core/src/mal\_core/ingest/daily\_nc.py:8--11}.

The series must contain every day of the rainfall axis. The check is explicit, the error message
names the first absent day, and the remedy names the download command
\src{mal-core/src/mal\_core/ingest/daily\_nc.py:270--278}. The re-indexing is done by an
explicit day lookup rather than by positional slicing, so a non-contiguous or misaligned
temperature axis does not silently shift the series
\src{mal-core/src/mal\_core/ingest/daily\_nc.py:279--284}.

Open-sea cells receive a temperature of exactly zero, not a missing value, and the reason is
recorded in the code: zero is the convention the engine's climate layer and the renderer already
understand, whereas a missing value would be imputed from neighbouring land cells at load time
and would quietly turn sea pixels into plausible land temperatures
\src{mal-core/src/mal\_core/ingest/daily\_nc.py:285--293}. The engine's reader converts the
sentinel to a genuine missing value on load
\src{mal-core/src/mal\_core/abm/src/env\_reader.cpp:207--216}, and the climate layer then
imputes it once, across every day \src{mal-core/src/mal\_core/abm/src/climate.cpp:167--177}. The
two behaviours are complementary: the \emph{builder} writes a zero, and the \emph{reader}
refuses to treat a large negative value as a temperature.

\subsection{The time span, and why the whole series is kept}

The NetCDF covers the whole daily series available at build time, not only the simulated window,
and the file name carries the first and last year of the rainfall axis
\src{mal-core/src/mal\_core/ingest/daily\_nc.py:465--466}. The justification is stated in the
module's own preamble: the days before the simulation start are the causal warm-up window over
which the engine integrates the pluvial state
\src{mal-core/src/mal\_core/ingest/daily\_nc.py:33--40}. The engine performs that integration
over the bands that precede the simulation start
\src{mal-core/src/mal\_core/abm/src/coordinator.cpp:73--90}. A consequence follows: the number
of days in the file is an input to the warm-up, and truncating the file would change the initial
state of every cell.

\section{The static planes, the coastal band, and the ocean mask} \label{sec:statics}

\subsection{How the static planes are assembled}

The \codeconst{ingest.env_static_planes} static planes are read: the permanent-water
fraction from the JRC product, and four
land-cover areal fractions from Dynamic World
\src{mal-core/src/mal\_core/ingest/daily\_nc.py:317--325}. A plane that is missing from its file
raises an error listing the variables actually present
\src{mal-core/src/mal\_core/ingest/daily\_nc.py:300--304}; a plane whose shape is not the
environment grid raises as well \src{mal-core/src/mal\_core/ingest/daily\_nc.py:310--313}.

A three-dimensional variable is a multi-year stack, and it is averaged over its year axis into a
single plane \src{mal-core/src/mal\_core/ingest/daily\_nc.py:305--309}. The Dynamic World loader
writes one frame per year, each being that year's modal-label composite reduced to areal
fractions \src{mal-commonlib/src/mal\_commonlib/data/loaders/dynamic\_world.py:336--409}; the
ingest builder averages those per-year planes. The consequence is that a multi-year run's static
potential is the mean of the years' fractions, not the fraction of the modal year --- a
deliberate smoothing, and a named choice rather than a default.

NaN in a static plane becomes zero \src{mal-core/src/mal\_core/ingest/daily\_nc.py:314}. The
rationale is recorded in the habitat module: absence of land-cover evidence is not a wet signal
\src{mal-core/src/mal\_core/habitat/potential.py:31--34}.

\subsection{The open-sea rule, stated as a formula}

A cell is open sea when it is water by land fraction \emph{and} farther than the coastal band
from land:

\begin{equation}
\text{open sea}(i) \;\iff\; \text{land\_frac}(i) < \text{sea\_water\_floor} \;\wedge\;
\text{dist\_to\_land}(i) > \text{coast\_band\_m}
\end{equation}

\noindent with \text{sea\_water\_floor} $= \mskey{ingest.sea_water_floor}$ and
\text{coast\_band\_m} $= \mskey{ingest.coastline_buffer_m}$ metres by default
\src{mal-core/src/mal\_core/ingest/daily\_nc.py:57--69}. The equation is the rule for a
\emph{positive} band; the zero band is a different branch of the loader's \code{sea\_mask} and
is stated below. The band is a \emph{distance query},
not a mask dilation: the code records that a former dilation moved the coast kilometres inland
\src{mal-core/src/mal\_core/ingest/daily\_nc.py:64--67}. The band width can be overridden for an
experiment through an environment variable, and a negative value disables the filter altogether
\src{mal-core/src/mal\_core/ingest/daily\_nc.py:240--250}.

A band of zero is not the equation above read with \text{coast\_band\_m} $= 0$: the loader
returns the water-only mask when the band is zero or negative, so only the land fraction
decides \src{mal-commonlib/src/mal\_commonlib/data/loaders/coastline.py:292--294}. A
\emph{negative} band disables the filter outright instead, because the builder then computes no
mask at all \src{mal-core/src/mal\_core/ingest/daily\_nc.py:249--253}: the rainfall, the
temperature and the \codeconst{ingest.env_static_planes} static planes keep their open-ocean
values.

\subsection{What the rule does, and why}

On open sea, the stage sets the rainfall to zero, the temperature to zero, and \emph{all} the
\codeconst{ingest.env_static_planes} static planes to zero \src{mal-core/src/mal\_core/ingest/daily\_nc.py:252--253},
\src{mal-core/src/mal\_core/ingest/daily\_nc.py:291--293},
\src{mal-core/src/mal\_core/ingest/daily\_nc.py:340--347}. The reason is recorded at the point
of the code that does it: Dynamic World reports its water class at one over the ocean, so
without the filter every ocean cell would carry a positive total potential and would enter the
candidate set \src{mal-core/src/mal\_core/ingest/daily\_nc.py:327--339}.

The coastal band keeps its statics. Lagoons, estuaries and river mouths are real breeding
habitat for these species, and the band is where the satellite imagery shows the shallow plume
\src{mal-core/src/mal\_core/ingest/daily\_nc.py:330--332}. So the choice is a named one: a
distance threshold on the open-ocean class, with the near-shore band deliberately retained.

\subsection{The censuses, and what they make auditable}

Two one-line censuses are printed. The saltwater filter line reports the before and after sums
of the permanent-water plane and of the open-plus-water plane, together with the number of cells
kept \src{mal-core/src/mal\_core/ingest/daily\_nc.py:349--356}. The coastline census reports the
number of cells that are land (more than half land), mixed, coastal water, and open sea
\src{mal-core/src/mal\_core/ingest/daily\_nc.py:364--374}. Together they let a reader confirm,
from the log alone, what the mask did.

\subsection{The optional salinity channel}

When the salinity slot resolves, a monthly sea-surface-salinity field is reprojected onto the
environment grid and each month is copied onto every day of that month
\src{mal-core/src/mal\_core/ingest/daily\_nc.py:389--401}. Land and no-data cells stay missing
in the source and become zero, the freshwater value
\src{mal-core/src/mal\_core/ingest/daily\_nc.py:550--557},
\src{mal-core/src/mal\_core/ingest/daily\_nc.py:401}. The variable is written with units of
practical salinity units \src{mal-core/src/mal\_core/ingest/daily\_nc.py:452--457}.

The channel is optional in the contract in both directions: its absence is normal, and the
engine reads an environment NetCDF without it as ``no salinity data'' and treats the cells as
freshwater \src{mal-core/src/mal\_core/abm/src/env\_reader.cpp:254--272}. When it is present it
acts as a hard gate: a patch whose salinity reaches the species' high tolerance cannot breed,
regardless of its accuracy \src{mal-core/src/mal\_core/abm/src/coordinator.cpp:221--225}.

\section{The habitat model} \label{sec:habitat}

\subsection{Six habitat types, two temporal classes}

The habitat types are frozen in the shared constants file, together with the temporal class of
each \src{mal-core/config/model\_constants.yaml:12--21}. The ordered list is
\mskey{habitat_types.types}, and the temporal class of each is a named constant:
\mskey{habitat_types.dynamics.surface_persistent},
\mskey{habitat_types.dynamics.wetland_persistent} and
\mskey{habitat_types.dynamics.urban_persistent} are
\mskey{habitat_types.dynamics.surface_persistent}-class persistent habitats;
\mskey{habitat_types.dynamics.wetland_pluvial}, \mskey{habitat_types.dynamics.open_pluvial} and
\mskey{habitat_types.dynamics.crop_pluvial} are pluvial.

The validator enforces two invariants on this list. The dynamics keys must be exactly the type
list \src{mal-core/src/mal\_core/config\_validation.py:379--382}, and every species' affinity
row must carry a weight for exactly those types and sum to one
\src{mal-core/src/mal\_core/config\_validation.py:395--416}. A third invariant is specific: the
\code{wetland\_pluvial} weight must be zero in every row, because the engine has no zone for
that type and a non-zero weight would be silently dropped on the C++ side
\src{mal-core/src/mal\_core/config\_validation.py:417--423}.

\subsection{The static potential $H[i,g]$}

The static potential is a pure function of the \codeconst{ingest.env_static_planes} feature
planes, computed in one place
\src{mal-core/src/mal\_core/habitat/potential.py:50--80}. Writing $P$ for the permanent-water
fraction, $V$ for the flooded-vegetation fraction, $U$ for the built fraction, $O$ for the
open-plus-water fraction and $C$ for the crop fraction, the six values are

\begin{align}
H_{\text{surface\_persistent}} &= P, & H_{\text{wetland\_persistent}} &= V, &
H_{\text{urban\_persistent}} &= \beta_{\text{urban}}\, U, \\ H_{\text{wetland\_pluvial}} &= V, &
H_{\text{open\_pluvial}} &= O + \delta_{\text{built}}\, U, & H_{\text{crop\_pluvial}} &= C,
\end{align}

\noindent where $\beta_{\text{urban}} = \mskey{urban.beta_urban}$ and $\delta_{\text{built}} =
\mskey{urban.rain_pools_built_damage}$ \src{mal-core/src/mal\_core/habitat/potential.py:44--49}.
Every value is clipped to $[0,1]$, and a missing feature is an error rather than a zero
\src{mal-core/src/mal\_core/habitat/potential.py:96--103}. The total is the sum of the six,
clipped again to $[0,1]$ \src{mal-core/src/mal\_core/habitat/potential.py:82--90}.

Two design decisions are recorded in the module's own preamble and are worth repeating. The
urban persistent contribution is not gated by the permanent-water fraction or by rainfall: a
cell with built cover and no recorded water keeps a small persistent offer through the dry
season \src{mal-core/src/mal\_core/habitat/potential.py:29--31}. And there is no weighting
matrix: the earlier linear-weight design was dropped, and resurrecting it would require a new
decision \src{mal-core/src/mal\_core/habitat/potential.py:34--36}.

\subsection{What the numbers are not} \label{sec:hnot}

$H$ is an \emph{index of relative offer} in $[0,1]$. It is not a measured pool count, not a
water depth, not a water volume, and not a probability that a given cell contains a given number
of larvae. The habitat module says so in its own preamble
\src{mal-core/src/mal\_core/habitat/potential.py:6--12}, and the contract table repeats it. The
consequences for a reader are:

\begin{itemize}
\item $H$ is dimensionless and has no units. It can be compared between cells of the same grid
and between runs of the same AOI; it cannot be read as ``this cell has this much water''. 
\item
The total is a \emph{sum over types}, so it is only meaningful as an index of overall offer. It
is used as such: as the candidate filter and as the oviposition-attraction amplitude. 
\item The
urban and open-pluvial contributions are scaled by two engineering constants, so the index is a
construct, not a measurement. Their provenance is declared in \S\ref{sec:provenance}.
\end{itemize}

\subsection{The two temporal responses are computed by the engine, not here} \label{sec:dynamic}

The habitat model has two shared temporal responses $A[i,k,t]$, indexed by the temporal class
$k$ of the habitat type:

\begin{itemize}
\item \textbf{persistent}: $A$ is the constant \mskey{capacity.persistent_availability} on every
day \src{mal-core/src/mal\_core/habitat/availability.py:5--9}. The persistent offer is
deliberately insensitive to drought; that is a documented minimum-viable-product limitation, not
an oversight \src{mal-core/src/mal\_core/habitat/availability.py:5--9}. 
\item \textbf{pluvial}:
$A$ is a state $z[i,t] \in [0,1]$ integrated day by day from the rainfall channel this stage
writes,
    \begin{equation}
z_{\text{dry}} = z \,\exp\!\left(-\Delta t / \tau_{\text{dry}}\right), \qquad z = z_{\text{dry}}
+ (1 - z_{\text{dry}})\bigl(1 - \exp(-r / r_{\text{scale}})\bigr),
    \end{equation}
with $\tau_{\text{dry}} = \mskey{pluvial.dry_memory_days}$, $r_{\text{scale}} =
\mskey{pluvial.rain_scale_mm}$ and $\Delta t = \mskey{pluvial.delta_days}$
\src{mal-core/src/mal\_core/habitat/availability.py:38--56}. The C++ mirror is the same
recursion
\src{mal-core/src/mal\_core/abm/include/mal\_abm\_fast/habitat\_availability.hpp:98--113}, and
the engine calls it once per cell per day
\src{mal-core/src/mal\_core/abm/src/coordinator.cpp:119--125}.
\end{itemize}

The Python availability module is a reference mirror, and it has \emph{no production consumer
at all}: the only modules that import it are two test modules
\src{mal-core/tests/test\_scoring\_pools.py:25},
\src{mal-core/tests/test\_habitat\_python\_contract.py:19--25}. The calibration scorer that
exercises the pluvial recursion does not import it either --- it imports the capacity module
and re-implements the $z$ step in place
\src{mal-core/src/mal\_core/scoring/scorers/d19\_pool\_persistence.py:23--24},
\src{mal-core/src/mal\_core/scoring/scorers/d19\_pool\_persistence.py:66--72}. The live
integration is the C++ code in the coordinator
\src{mal-core/src/mal\_core/habitat/availability.py:1--30}. The consequence for this document is
worth stating plainly: \emph{the ingest stage writes the forcing for $A$; it does not write $A$
itself}.

\subsection{$S$ and $K$ are computed by the engine, not here}

The daily accuracy and the daily capacity are engine quantities:

\begin{equation}
S[i,s,t] = \operatorname{clip}\Bigl(\sum_g w[s,g]\, H[i,g]\, A[i,k(g),t],\; 0,\; 1\Bigr), \qquad
K[i,s,t] = k_{\text{ref}} \cdot \text{cell\_area\_km}^2[i] \cdot \max(0, S[i,s,t]).
\end{equation}

\noindent The engine computes them per cell per day
\src{mal-core/src/mal\_core/abm/src/coordinator.cpp:129--149},
\src{mal-core/src/mal\_core/abm/include/mal\_abm\_fast/habitat\_availability.hpp:159--185}, with
$k_{\text{ref}} = \mskey{species.k_ref}$, an absolute per-cell capacity in individuals per day
per square kilometre \src{mal-core/config/model\_constants.yaml:47--52}.

The Python reference module computes $S$ and $K$ for the scorer path with a \emph{different}
capacity constant, a diagnostic one, explicitly documented as not being the engine's:
\codeconst{ingest.k_ref_diag_per_km2} individuals per square kilometre per day
\src{mal-core/src/mal\_core/habitat/capacity.py:26--30}. The two are different concepts and must
not be read as one number; the tombstone module that used to hold shared constants repeats the
distinction \src{mal-core/src/mal\_core/habitat/constants.py:1--16}. The reader who wants the
capacity that the simulation actually uses should take $k_{\text{ref}} = \mskey{species.k_ref}$
and the formula above (\S\ref{sec:open} question 9).

\subsection{The Python and C++ potentials do not compute the same set}

The two implementations agree on the five living values but not on the \emph{set} of types.

\begin{itemize}
\item Python produces six habitat types, including \code{wetland\_pluvial}, whose value is the
flooded-vegetation fraction again \src{mal-core/src/mal\_core/habitat/potential.py:71--79}. Its
total therefore adds the flooded-vegetation fraction a second time. 
\item The engine produces
six \emph{zones}, one of which is a permanently zero forest zone, and has no
\code{wetland\_pluvial} zone at all
\src{mal-core/src/mal\_core/abm/include/mal\_abm\_fast/habitat\_availability.hpp:56--85},
\src{mal-core/src/mal\_core/abm/include/mal\_abm\_fast/zone\_id.hpp:13--21}. Its total does not
include the second flooded-vegetation term
\src{mal-core/src/mal\_core/abm/include/mal\_abm\_fast/habitat\_availability.hpp:87--96}.
\end{itemize}

The consequence is a bounded but real difference between the two totals: the Python total
exceeds the engine's by at most the flooded-vegetation fraction, which is at most one. The
quantity $S$ is unaffected, because the \code{wetland\_pluvial} affinity is zero in every
species row and the validator enforces it
\src{mal-core/src/mal\_core/config\_validation.py:417--423}. What is affected is the amplitude
written into the GeoPackage and the count printed by the ingest census --- both of which use the
Python total (\S\ref{sec:open} question 2).

\subsection{Three candidate floors, with three different roles}

\begin{center}
\small
\begin{longtable}{@{}p{3.3cm}p{2.4cm}p{8.0cm}@{}}
\toprule
Floor & Value & Role \\
\midrule
\endhead
Ingest census & \mskey{ingest.candidate_min_ingest}
  & The threshold above which a cell is counted as a candidate in the
environment build; an empty count is a hard error
\src{mal-core/src/mal\_core/ingest/daily\_nc.py:376--387}. \\ GeoPackage membership & $0$
  & The threshold above which a cell is written into the GeoPackage as a
point feature \src{mal-core/src/mal\_core/ingest/env.py:70}. \\ Engine floor &
\mskey{habitat.candidate_min_engine}
  & The threshold below which the engine drops a patch from the GeoPackage,
and the threshold at which the engine admits a cell into its own candidate union
\src{mal-core/src/mal\_core/abm/src/habitat\_engine.cpp:214--219},
\src{mal-core/src/mal\_core/abm/src/coordinator.cpp:169--175}. \\
\bottomrule
\end{longtable}
\end{center}

\smallskip \noindent The constants file states that the ingest floor and the engine floor have
different roles and are deliberately not mirrors of each other
\src{mal-core/config/model\_constants.yaml:151--153}. The three thresholds are therefore not
three copies of one number, and a reader comparing the ingest census with the engine's candidate
count must not expect them to agree (\S\ref{sec:open} question 10).

\section{The habitat GeoPackage, field by field} \label{sec:gpkg}

\subsection{What is written}

The GeoPackage holds one Point feature per candidate cell, positioned at the cell centre in the
AOI's coordinate reference system \src{mal-core/src/mal\_core/ingest/env.py:69--73},
\src{mal-core/src/mal\_core/ingest/env.py:99--103}. The columns are:

\begin{center}
\small
\begin{longtable}{@{}p{3.5cm}p{9.4cm}@{}}
\toprule
Column & Meaning \\
\midrule
\endhead
\code{hab\_type} & The dominant habitat type: the argument maximum of the six potential values
at that cell \src{mal-core/src/mal\_core/ingest/env.py:81--85}. \\ \code{dynamics} & The
temporal class of that dominant type, written as the string \code{persistent} or \code{pluvial}
\src{mal-core/src/mal\_core/ingest/env.py:86--90}. \\ \code{habitat\_potential} & The clipped
total potential of the cell, the amplitude used for oviposition attraction
\src{mal-core/src/mal\_core/ingest/env.py:91}. \\ \code{row}, \code{col} & The grid identity of
the cell, the pair the engine indexes by \src{mal-core/src/mal\_core/ingest/env.py:92--93}. \\
\code{aoi\_slug} & The slug of the area the file belongs to
\src{mal-core/src/mal\_core/ingest/env.py:94}. \\ \code{hab\_<type>} & One column per habitat
type, carrying that type's per-cell potential \src{mal-core/src/mal\_core/ingest/env.py:96--97}.
\\
\end{longtable}
\end{center}

\smallskip \noindent There are deliberately no empty beacon columns: the retired
\code{twi\_value} and \code{urban\_capacity\_factor} columns are gone
\src{mal-core/src/mal\_core/ingest/env.py:60--62}.

\subsection{What the engine actually reads}

The engine resolves the fields \code{hab\_type}, \code{habitat\_potential}, \code{dynamics},
\code{row} and \code{col} \src{mal-core/src/mal\_core/abm/src/habitat\_engine.cpp:165--171}. It
accepts any feature regardless of \code{hab\_type} --- viability is decided on potential, not on
the label \src{mal-core/src/mal\_core/abm/src/habitat\_engine.cpp:190--192}. It does \emph{not}
read the \code{hab\_<type>} columns: it recomputes the potential from the environment NetCDF's
static planes \src{mal-core/src/mal\_core/abm/src/climate.cpp:196--215},
\src{mal-core/src/mal\_core/abm/include/mal\_abm\_fast/habitat\_availability.hpp:56--85}. The
Python-side potential in the GeoPackage and the C++-side potential are therefore two
implementations of the same idea, differing as described in \S\ref{sec:habitat}. The module
preamble's claim that the C++ parser requires the H columns
\src{mal-core/src/mal\_core/ingest/env.py:60--62} is contradicted by the reader, and the reader
wins.

\subsection{Verified on the delivered Ghana file}

Reading \code{data/ghana/ghana\_habitat\_patches.gpkg} with a geospatial reader gives 382\,127
Point features on the 779 by 551 Ghana grid --- 89\,\% of the cells. The \code{hab\_type}
distribution is \code{open\_pluvial} 326\,608, \code{urban\_persistent} 26\,696,
\code{surface\_persistent} 15\,069, \code{crop\_pluvial} 13\,316 and \code{wetland\_persistent}
438; the \code{dynamics} column splits 339\,924 \code{pluvial} and 42\,203 \code{persistent}.
The amplitude ranges from $3\times10^{-6}$ to one, with a median near 0.725. The count agrees
with the \code{n\_candidate\_cells} attribute of the environment NetCDF, so the census and the
GeoPackage agree in that run.

The dominance of \code{open\_pluvial} is a direct consequence of the Dynamic World mapping: the
open-plus-water plane includes the water class, so nearly every non-ocean cell carries some
open-pluvial potential.

\subsection{The delivered file was produced by the pre-fix implementation}

The delivered artefacts do not satisfy the open-sea rule this document describes, because they
were produced by an earlier implementation of the builder. The delivered environment NetCDF
carries no \code{land\_frac} plane, although the builder writes one whenever the coastline
slot resolves \src{mal-core/src/mal\_core/ingest/daily\_nc.py:443--447}; and the stored
coastline artefact \code{data/ghana/ghana\_land\_mask.tif} is a binary $\{0,1\}$ mask, whereas
the current loader returns a per-cell land fraction
\src{mal-commonlib/src/mal\_commonlib/data/loaders/coastline.py:297--330}. Applying the
current rule to the stored statics marks 45\,674 cells open sea, all of them present in the
GeoPackage, and would reduce the candidate count from 382\,127 to 336\,453. The feature count,
the type distribution and the amplitudes verified above are therefore that earlier
implementation's, not the current builder's (\S\ref{sec:open} question 19).

\subsection{Two disclosed defects}

\begin{description}
\item[No version on the file.] The GeoPackage carries neither a contract version nor a generator
version. The writer writes no metadata \src{mal-core/src/mal\_core/ingest/env.py:99--105}, and
the delivered file has no \code{gpkg\_metadata} table at all. Only the environment NetCDF is
versioned (\S\ref{sec:envnc}). A reader of the GeoPackage has no way to tell which mapping
produced it (\S\ref{sec:open} question 3). 
\item[The \code{dynamics} column is a string and is
read as an integer.] The writer emits \code{persistent} or \code{pluvial}
\src{mal-core/src/mal\_core/ingest/env.py:86--90}; the reader fetches the field with an integer
accessor \src{mal-core/src/mal\_core/abm/src/habitat\_engine.cpp:220}, which on a string field
yields zero. The engine's per-patch dynamics flag is therefore always zero, including for
pluvial patches. The field is diagnostic only, so the simulation is unaffected and the
diagnostic is wrong (\S\ref{sec:open} question 4).
\end{description}

\section{The environment NetCDF contract} \label{sec:envnc}

\subsection{Variables, dimensions, attributes, encoding}

\begin{center}
\small
\begin{longtable}{@{}p{3.4cm}p{2.6cm}p{2.4cm}p{8.0cm}@{}}
\toprule
Variable & Dimensions & Units & Role \\
\midrule
\endhead
\code{rainfall} & (time, y, x) & mm/day & The daily precipitation forcing \\ \code{air\_temp\_c}
& (time, y, x) & degC & The daily thermal forcing \\ \code{jrc\_persistent} & (y, x) & 1 &
Permanent-surface-water fraction \\ \code{lc\_built} & (y, x) & 1 & Built areal fraction \\
\code{lc\_wetland} & (y, x) & 1 & Flooded-vegetation areal fraction \\
\code{lc\_open\_plus\_water} & (y, x) & 1 & Open-plus-water areal fraction \\ \code{lc\_crop} &
(y, x) & 1 & Crop areal fraction \\ \code{land\_frac} & (y, x) & 1 & Optional display basemap;
not part of the habitat feature bundle \\ \code{salinity\_ppt} & (time, y, x) & psu & Optional
monthly-broadcast salinity \\
\end{longtable}
\end{center}

\smallskip \noindent The \codeconst{ingest.env_static_planes} static planes are
two-dimensional and are never broadcast onto the time axis; the
\codeconst{ingest.env_daily_channels} daily channels never enter the static planes
\src{mal-core/src/mal\_core/ingest/daily\_nc.py:404--431}. The global attributes are the
conventions string, the title, the AOI slug, the scale, the contract version, the generator
version and the candidate count \src{mal-core/src/mal\_core/ingest/daily\_nc.py:433--441}. Every
variable is written as 32-bit float with deflate compression at the configured level
\src{mal-core/src/mal\_core/ingest/daily\_nc.py:459--463}.

\subsection{The contract version}

The contract version is the generated constant \mskey{contracts.env_nc}, read from the shared
constants file \src{mal-core/src/mal\_core/ingest/daily\_nc.py:53} and written into the
attribute \src{mal-core/src/mal\_core/ingest/daily\_nc.py:438}. The validator requires the value
to have the form \emph{major.minor} \src{mal-core/src/mal\_core/config\_validation.py:368--373}.
The state contract, \mskey{contracts.cogs_state}, is a different contract belonging to the
simulation stage and is not written by this stage.

The generator version is \emph{not} generated. It is a literal string in the builder
\src{mal-core/src/mal\_core/ingest/daily\_nc.py:439} and has no key in the constants file. This
document therefore quotes it as a reading of the code rather than through the normal mechanism
(\S\ref{sec:open} question 8).

\subsection{NoData convention}

The pipeline-wide NoData sentinel is \mskey{ingest.nodata_sentinel}
\src{mal-core/src/mal\_core/ingest/daily\_nc.py:55}. The environment NetCDF does not write an
explicit fill value: the encoding block sets the data type and the compression only
\src{mal-core/src/mal\_core/ingest/daily\_nc.py:459--463}, and what lands in the file is the
writer's default NaN fill for 32-bit floats. The delivered Ghana file confirms it: every
variable carries \code{\_FillValue = NaNf}. The sentinel is used \emph{during} the build, to
detect rainfall NoData \src{mal-core/src/mal\_core/ingest/daily\_nc.py:209}, and is not stored.

\subsection{Verified on the delivered Ghana file}

The file \code{data/ghana/ghana\_regional\_2024\_2024\_env.nc} has the two daily channels with
shape 366 by 779 by 551 in 32-bit float, the \codeconst{ingest.env_static_planes} static
planes with shape 779 by 551 in 32-bit float, and exactly these global attributes: \code{Conventions: CF-1.8}, a title,
\code{aoi\_slug: ghana}, \code{scale: regional}, \code{contract\_version: 3.0},
\code{generator\_version: proxy-habitat-v1-0.1.0} and \code{n\_candidate\_cells: 382127}. It
carries no \code{land\_frac} plane and no \code{salinity\_ppt} plane, and no coordinate
reference system (\S\ref{sec:frame}). The absence of the \code{land\_frac} plane is not a
property of the file's contents that the builder would choose: the builder writes that plane
whenever the coastline slot resolves
\src{mal-core/src/mal\_core/ingest/daily\_nc.py:443--447}. Its absence shows that the
delivered file was produced by the pre-fix implementation, together with the habitat
GeoPackage (\S\ref{sec:gpkg}, \S\ref{sec:open} question 18).

\section{Hosts, buildings and the wildlife proxy} \label{sec:hosts}

\subsection{The nine planes}

The host grid is a two-dimensional NetCDF on the same AOI grid, with one plane per host layer
\src{mal-core/src/mal\_core/ingest/hosts.py:186--213}. The layers are: human population; five
livestock species --- cattle, goats, sheep, pigs and chickens
\src{mal-core/src/mal\_core/ingest/hosts.py:33}; the urban class; the building fraction; and the
wildlife host proxy. A companion JSON file carries per-variable statistics and the grid
specification \src{mal-core/src/mal\_core/ingest/hosts.py:208--214}.

\subsection{The aggregation rules, and what each one loses}

\begin{center}
\small
\begin{longtable}{@{}p{3.0cm}p{2.2cm}p{9.3cm}@{}}
\toprule
Layer & Rule & What the rule means and what it costs \\
\midrule
\endhead
Human population & sum
  & A count-preserving aggregation onto the AOI grid
\src{mal-core/src/mal\_core/ingest/hosts.py:75--88}. The number in a cell is a population count.
\\ Livestock (five species) & sum
  & The source product is a \emph{density} in heads per square kilometre
\src{mal-commonlib/src/mal\_commonlib/data/loaders/glw.py:13--16},
\src{mal-commonlib/src/mal\_commonlib/data/loaders/glw.py:224--227}; summing densities yields a
sum of densities, \emph{not} a head count. The engine consumes these numbers as counts, which is
a stated inconsistency worth carrying in mind. \\ Urban class & nearest
  & A categorical layer, so it is resampled by nearest neighbour
\src{mal-core/src/mal\_core/ingest/hosts.py:112--120}, then reclassified:
\codeconst{ingest.urban_class_urban} stays \codeconst{ingest.urban_class_urban} (urban),
twenty to twenty-nine becomes \codeconst{ingest.urban_class_rural} (rural), zero to nineteen
becomes \codeconst{ingest.urban_class_water} (water)
\src{mal-core/src/mal\_core/ingest/hosts.py:122--128}. The reclassification is \emph{inverted}
with respect to the product it reclassifies (\S\ref{sec:urbanclass}). \\ Building fraction & sum
  & The source is a per-cell presence of building footprint centroids, at
one where any centroid falls in the cell
\src{mal-commonlib/src/mal\_commonlib/data/loaders/buildings.py:187--189}; the quantity is
therefore a centroid-presence indicator, not an area fraction. \\ Wildlife host proxy & nearest
  & A categorical suitability surface, resampled by nearest neighbour
\src{mal-core/src/mal\_core/ingest/hosts.py:170--178}. \\
\end{longtable}
\end{center}

\subsection{The urban class is inverted with respect to the product} \label{sec:urbanclass}

The reclassification above does not match the GHS-SMOD product it reclassifies. In the R2023A
scheme the codes are 30 for an urban centre, 21, 22 and 23 for urban clusters of low, medium
and high density, 11, 12 and 13 for rural low, medium and high density, and 10 for water; the
loader's own comment reads the scheme the other way round, calling 20--29 rural
\src{mal-commonlib/src/mal\_commonlib/data/loaders/ghsl.py:76--79}, and the reclassification
follows that reading. The consequence is that the product's urban clusters become ``rural''
(fifty) and its rural classes become ``water'' (twenty); only the urban centre (thirty) and
the water class (ten) keep the right label.

The delivered Ghana host grid shows the inversion. Its \code{urban\_class} plane carries
415\,992 cells labelled twenty, 9\,434 labelled fifty and 3\,803 labelled thirty, while the
source product's own classes on the same grid are 30: 3\,803 cells, 23: 1\,990, 22: 594,
21: 6\,850, 13: 7\,457, 12: 51\,101, 11: 299\,511, 10: 57\,490 and 0: 433. The engine maps
the class to an indoor fraction with a threshold of forty, so that anything at or below forty
is urban \src{mal-core/src/mal\_core/abm/src/host\_landscape.cpp:111--118}. Every one of the
415\,992 cells labelled twenty is therefore read by the engine as urban: they are the product's
rural cells (358\,069 of them) together with its water cells, and they make up roughly
97\,\% of the area of interest, while the 9\,434 genuine urban clusters are read as rural. The
urban/rural signal is inverted over the area of interest (\S\ref{sec:open} question 17).

\subsection{The three silent-ish degradations}

If the urban-class loader fails, every cell becomes the rural default
\codeconst{ingest.urban_class_rural}
\src{mal-core/src/mal\_core/ingest/hosts.py:133--136}. If the building loader fails, the
fraction becomes \codeconst{ingest.fallback_building_urban} in urban cells and
\codeconst{ingest.fallback_building_rural} elsewhere
\src{mal-core/src/mal\_core/ingest/hosts.py:156--159}. If the wildlife loader fails, the proxy
becomes a constant \codeconst{ingest.fallback_wildlife_proxy} everywhere
\src{mal-core/src/mal\_core/ingest/hosts.py:180--182}. Each
degradation is announced on the standard error stream, and each substitutes a constant for a
real layer. The three constants are literals in the module rather than constants-file keys, so
this document quotes them through the generated code-constant macros, which read the module
itself (\S\ref{sec:open} question 8).

The two builder flags make the same substitutions deliberately: skipping the buildings layer
uses the urban heuristic \src{mal-core/src/mal\_core/ingest/hosts.py:141}, and skipping the
wildlife layer uses the constant \src{mal-core/src/mal\_core/ingest/hosts.py:164}.

\section{Manifest integration} \label{sec:manifest}

\subsection{How the stage writes to the manifest}

Every artefact is registered through one shared helper
\src{mal-core/src/mal\_core/ingest/\_shared.py:7--33}, which delegates to the manifest writer
\src{mal-core/src/mal\_core/download/manifest.py:118--189}. The entry records the type, the
format, the \code{required\_for\_abm} flag, the file name, an optional variable list, an
optional period, a \code{+}-joined compound licence identifier and a literal attribution line
\src{mal-core/src/mal\_core/download/manifest.py:151--180}. The manifest's
\code{expected\_files} list is rebuilt from scratch on every write, as the union of every
declared file of every entry \src{mal-core/src/mal\_core/download/manifest.py:181--184}.

There are four registrations: two from the environment builder
\src{mal-core/src/mal\_core/ingest/env.py:147--162},
\src{mal-core/src/mal\_core/ingest/env.py:174--187} and two from the host builder
\src{mal-core/src/mal\_core/ingest/hosts.py:219--240}.

\subsection{The year key, and what it means for a multi-year file}

The environment NetCDF physically spans several years --- its name carries the first and last
year of the rainfall axis --- but it is registered under a single year key, the year the CLI was
given \src{mal-core/src/mal\_core/ingest/daily\_nc.py:465--466},
\src{mal-core/src/mal\_core/ingest/env.py:148}. The manifest writer stores a year-keyed file
under \code{files[str(year)]} and an unkeyed one under \code{files[dataset\_name]}
\src{mal-core/src/mal\_core/download/manifest.py:176--179}.

The consequence, as a described situation: the simulation wrapper resolves the environment
artefact \emph{with} a year \src{mal-core/src/mal\_core/abm/wrapper.py:137--141}. If the ingest
run was told about one year and the simulation is asked for another, the resolution fails ---
even though the file on disk contains the requested year's days. If a manifest entry is per-year
keyed and lacks the requested year, the resolution raises rather than silently picking another
year \src{mal-core/src/mal\_core/download/manifest.py:91--98}. The habitat GeoPackage is
registered under a year key too, but is resolved without one, so the fallback to the first
declared file applies \src{mal-core/src/mal\_core/abm/wrapper.py:142--146} (\S\ref{sec:open}
question 7).

\subsection{The declared variable list does not match the file}

The environment NetCDF's registered variable list is derived from the built file's own variable
list \src{mal-core/src/mal\_core/ingest/daily\_nc.py:498--503}, and the delivered manifest
matches the delivered file exactly. The habitat GeoPackage's registered list does not. It is
constructed from the \emph{feature plane} names
\src{mal-core/src/mal\_core/ingest/env.py:183--185}, while the file itself carries one
\code{hab\_<habitat type>} column per habitat type
\src{mal-core/src/mal\_core/ingest/env.py:96--97}, plus the \code{dynamics} and \code{aoi\_slug}
columns, which the registration omits.

Verified on the delivered artefacts: the manifest advertises the habitat entry's variables as
\code{hab\_type}, \code{habitat\_potential}, \code{row}, \code{col},
\code{hab\_jrc\_persistent}, \code{hab\_lc\_built}, \code{hab\_lc\_wetland},
\code{hab\_lc\_open\_plus\_water} and \code{hab\_lc\_crop}, while the file carries
\code{hab\_surface\_persistent}, \code{hab\_wetland\_persistent}, \code{hab\_urban\_persistent},
\code{hab\_wetland\_pluvial}, \code{hab\_open\_pluvial}, \code{hab\_crop\_pluvial},
\code{dynamics} and \code{aoi\_slug}. Five advertised names are not in the file and eight present
names are not advertised (\S\ref{sec:open} question 5).

\section{The seam with the simulation} \label{sec:seam}

\subsection{What the engine resolves}

Before running, the wrapper requires that every file declared in the manifest exists
\src{mal-core/src/mal\_core/abm/wrapper.py:112--121}. It then resolves four paths: the
environment NetCDF with a year, the habitat GeoPackage, the host grid and the wind field
\src{mal-core/src/mal\_core/abm/wrapper.py}. The wind field is
\emph{not} produced by this stage --- it comes from the download stage directly
\src{mal-core/src/mal\_core/abm/wrapper.py:149--152}. The environment NetCDF and the habitat
GeoPackage are the only two hard requirements; the rest are passed to the engine as flags when
they resolve \src{mal-core/src/mal\_core/abm/wrapper.py:166--178},
\src{mal-core/src/mal\_core/abm/wrapper.py:210--222}.

\subsection{What the engine does with the environment NetCDF}

The engine reads the \codeconst{ingest.env_daily_channels} daily channels and the
\codeconst{ingest.env_static_planes} static planes
\src{mal-core/src/mal\_core/abm/src/env\_reader.cpp:204--246}. The optional \code{land\_frac}
plane is a display basemap and is explicitly not part of the habitat feature bundle
\src{mal-core/src/mal\_core/abm/src/climate.cpp:211--213}. The optional \code{salinity\_ppt}
plane's absence means ``no salinity data'': the reader simply leaves the salinity vector empty
\src{mal-core/src/mal\_core/abm/src/env\_reader.cpp:254--272}, and the accessor then returns
the configured freshwater anchor \mskey{aquatic_stages.salinity_missing_psu}
\src{mal-core/src/mal\_core/abm/src/climate.cpp:149--158}. A static plane
whose size does not match the grid is silently treated as zero, with a printed note
\src{mal-core/src/mal\_core/abm/src/climate.cpp:196--205}. When the static planes are absent
altogether, the engine prints its own complaint about the environment NetCDF's contract version
\src{mal-core/src/mal\_core/abm/src/coordinator.cpp:66--70}.

\subsection{How the engine forms its candidate set}

The engine's patch set is the \emph{union} of two sources: the patches read from the GeoPackage,
and every cell whose recomputed total potential passes the engine's floor
\src{mal-core/src/mal\_core/abm/src/coordinator.cpp:161--175}. The GeoPackage is thus a
conservative starting set rather than the authoritative one, and the two candidate counts need
not agree (\S\ref{sec:open} question 10).

Each patch's daily state is then computed from the recomputed potential, the current pluvial
state and the salinity channel \src{mal-core/src/mal\_core/abm/src/coordinator.cpp:207--226}. A
patch is activated when its accuracy reaches \mskey{habitat.activation_min_engine}, and is
deactivated regardless of accuracy when its salinity reaches the species' high tolerance.

\subsection{What the engine recomputes}

The engine recomputes the potential, the two temporal responses, the daily accuracy and the
daily capacity from the static planes and the rainfall channel, every day (\S\ref{sec:dynamic}).
It therefore consumes the ingest artefacts as \emph{static structure plus forcing}: the grid,
the candidate cells, the amplitudes, and the daily rainfall and temperature. Everything dynamic
about habitat is downstream.

\section{Declaration of absences} \label{sec:absences}

The following are \emph{not} done by this stage, and a reader looking for them should not find
them:

\begin{itemize}
\item \textbf{No terrain}: no digital elevation model, no topographic wetness index, no
catchment ratio, no slope, no drainage \src{mal-core/src/mal\_core/ingest/daily\_nc.py:36--40}.
\item \textbf{No millimetre water balance}: no pool depth, no storage volume, no evaporation
term, no per-cell aquatic solver. 
\item \textbf{No annual temperature-suitability surface}; the
retired channel and its broadcast path are gone
\src{mal-core/src/mal\_core/ingest/daily\_nc.py:8--11}. 
\item \textbf{No vegetation-index
channel} \src{mal-core/src/mal\_core/ingest/daily\_nc.py:36--40}. 
\item \textbf{No alpha or
theta weighting matrix} between potential and accuracy
\src{mal-core/src/mal\_core/habitat/potential.py:34--36}. 
\item \textbf{No species selection};
the potentials written are species-independent. 
\item \textbf{No capacity computation}: $S$ and
$K$ are engine quantities (\S\ref{sec:dynamic}). 
\item \textbf{No climatology, no trend
analysis, no gap interpolation} other than the rainfall donor fill
\src{mal-core/src/mal\_core/ingest/daily\_nc.py:206--223}. 
\item \textbf{No per-cell water-body
enumeration}; the representation is a cell-level index, not a catalogue of ponds. 
\item
\textbf{No wind product}; it comes from the download stage
\src{mal-core/src/mal\_core/abm/wrapper.py:149--152}.
\end{itemize}

\section{Provenance of every constant used here} \label{sec:provenance}

\begin{center}
\small
\begin{longtable}{@{}p{4.4cm}p{3.0cm}p{7.6cm}@{}}
\toprule
Constant & Tier & Basis \\
\midrule
\endhead
\mskey{contracts.env_nc} & model definition & The contract string the environment NetCDF
declares; validated for form \src{mal-core/src/mal\_core/config\_validation.py:368--373}. \\
\mskey{ingest.nodata_sentinel} & engineering prior & The pipeline-wide NoData value; frozen so
that the C++ and Python readers agree \src{mal-core/config/model\_constants.yaml:137}. \\
\mskey{ingest.coastline_buffer_m} & engineering prior & Width of the coastal band; a choice,
tunable per run by an environment variable
\src{mal-core/src/mal\_core/ingest/daily\_nc.py:240--242}. \\ \mskey{ingest.sea_water_floor} &
engineering prior & Land fraction below which a cell counts as water
\src{mal-core/config/model\_constants.yaml:148}. \\ \mskey{ingest.candidate_min_ingest} &
engineering prior & The census floor; a guard against an empty static field, not a biological
threshold \src{mal-core/config/model\_constants.yaml:151--153}. \\
\mskey{ingest.fill_frac_valid_max} & engineering prior & The NoData fraction above which a cell
may no longer donate \src{mal-core/config/model\_constants.yaml:155}. \\
\mskey{ingest.grid_degrees_step_max} & engineering prior & The condition under which CF latitude
and longitude variables are written \src{mal-core/config/model\_constants.yaml:159}. \\
\mskey{ingest.nc_compression_level} & engineering prior & File-format configuration, not a model
quantity \src{mal-core/config/model\_constants.yaml:160}. \\ \mskey{ingest.cell_size_hint_m} &
engineering prior & Fallback cell size when the rainfall file carries no resolution attribute
\src{mal-core/src/mal\_core/ingest/daily\_nc.py:243--247}. \\ \mskey{urban.beta_urban} &
engineering prior & Scales the built fraction into the urban persistent potential; no
measurement or literature anchor is recorded for it
\src{mal-core/config/model\_constants.yaml:43}. \\ \mskey{urban.rain_pools_built_damage} &
engineering prior & Adds a rain-pool contribution on the built surface to the open-pluvial
potential; likewise unanchored \src{mal-core/config/model\_constants.yaml:44}. \\
\mskey{capacity.persistent_availability} & model definition & The constant persistent response;
deliberately drought-insensitive \src{mal-core/config/model\_constants.yaml:129}. \\
\mskey{pluvial.rain_scale_mm} & engineering prior & The rainfall scale of the pluvial recharge
\src{mal-core/config/model\_constants.yaml:29}. \\ \mskey{pluvial.dry_memory_days} & engineering
prior & The drying memory of the pluvial state \src{mal-core/config/model\_constants.yaml:30}.
\\ \mskey{pluvial.delta_days} & model definition & The one-day integration step
\src{mal-core/config/model\_constants.yaml:31}. \\ \mskey{pluvial.cold_start_z} & model
definition & The state before the first simulated day
\src{mal-core/config/model\_constants.yaml:32}. \\ \mskey{species.k_ref} & engineering prior &
The engine's per-cell capacity scale \src{mal-core/config/model\_constants.yaml:51}. \\
\mskey{habitat.candidate_min_engine} & engineering prior & The engine's candidate floor
\src{mal-core/config/model\_constants.yaml:25}. \\ \mskey{habitat.activation_min_engine} &
engineering prior & The daily activation threshold
\src{mal-core/config/model\_constants.yaml:24}. \\
\end{longtable}
\end{center}

\smallskip \noindent Two tiers appear nowhere in this stage's own constants, and their absence
is worth noting: there is no \emph{measured} constant and no \emph{literature-anchored} constant
in the ingest block. Every number above is either a contract string or an engineering choice.
The species affinity rows, which are the closest thing to an anchored set, live per species and
are validated to sum to one \src{mal-core/src/mal\_core/config\_validation.py:413--416}; for the
default species row \mskey{species.default_species} the weights are
\mskey{species.coluzzii.weights.surface_persistent},
\mskey{species.coluzzii.weights.wetland_persistent},
\mskey{species.coluzzii.weights.urban_persistent},
\mskey{species.coluzzii.weights.wetland_pluvial}, \mskey{species.coluzzii.weights.open_pluvial}
and \mskey{species.coluzzii.weights.crop_pluvial}. Note that the coordinator that computes the
live accuracy uses the species-neutral row, not the default species row
\src{mal-core/src/mal\_core/abm/src/coordinator.cpp:117},
\src{mal-core/src/mal\_core/abm/include/mal\_abm\_fast/habitat\_availability.hpp:146--156}.

\section{The generated parameter and contract tables} \label{sec:contracts}

The constants of the namespaces this stage draws on, as generated from the shared constants
file, with the production files that act on each one:

% Generated by tools/gen_concepts_values.py -- do not edit.
\paragraph{ingest: model parameters}
\begin{longtable}{@{}llll@{}}
\toprule
Key & Value & Unit & Read by \\
\midrule
\endhead
\code{capacity.persistent\_availability} & 1.0 & adimensional & mal-core/src/mal\_core/abm/include/mal\_abm\_fast/config.hpp, mal-core/src/mal\_core/abm/include/mal\_abm\_fast/habitat\_availability.hpp \\
\code{contracts.cogs\_state} & "4.0" & string & mal-core/src/mal\_core/abm/include/mal\_abm\_fast/config.hpp, mal-core/src/mal\_core/abm/include/mal\_abm\_fast/output\_contract.hpp, mal-core/src/mal\_core/abm/include/mal\_abm\_fast/wire.hpp \\
\code{contracts.env\_nc} & "3.0" & string & (loaded, never used) \\
\code{geodesy.earth\_radius\_m} & 6371008.8 & m & mal-core/src/mal\_core/abm/include/mal\_abm\_fast/aoi.hpp, mal-core/src/mal\_core/abm/include/mal\_abm\_fast/config.hpp, mal-core/src/mal\_core/abm/include/mal\_abm\_fast/wire.hpp \\
\code{geodesy.metres\_per\_deg\_lat} & 111195.08023353292 & m/° & mal-core/src/mal\_core/abm/include/mal\_abm\_fast/aoi.hpp, mal-core/src/mal\_core/abm/include/mal\_abm\_fast/config.hpp … (+4 more) \\
\code{habitat\_types.dynamics.crop\_pluvial} & "pluvial" & categórico & mal-core/src/mal\_core/abm/src/config.cc \\
\code{habitat\_types.dynamics.open\_pluvial} & "pluvial" & categórico & mal-core/src/mal\_core/abm/src/config.cc \\
\code{habitat\_types.dynamics.surface\_persistent} & "persistent" & categórico & mal-core/src/mal\_core/abm/src/config.cc \\
\code{habitat\_types.dynamics.urban\_persistent} & "persistent" & categórico & mal-core/src/mal\_core/abm/src/config.cc \\
\code{habitat\_types.dynamics.wetland\_persistent} & "persistent" & categórico & mal-core/src/mal\_core/abm/src/config.cc \\
\code{habitat\_types.dynamics.wetland\_pluvial} & "pluvial" & categórico & mal-core/src/mal\_core/abm/include/mal\_abm\_fast/config.hpp, mal-core/src/mal\_core/abm/src/habitat\_engine.cpp \\
\code{habitat\_types.types} & "surface\_persistent, wetland\_persistent, urban\_persistent, wetland\_pluvial, open\_pluvial, crop\_pluvial" & --- & mal-core/src/mal\_core/abm/include/mal\_abm\_fast/bite\_ledger.hpp, mal-core/src/mal\_core/abm/include/mal\_abm\_fast/config.hpp … (+7 more) \\
\code{ingest.candidate\_min\_ingest} & 1e-06 & índice & mal-core/src/mal\_core/abm/include/mal\_abm\_fast/config.hpp, mal-core/src/mal\_core/ingest/daily\_nc.py \\
\code{ingest.cell\_size\_hint\_m} & 1000.0 & m & mal-core/src/mal\_core/abm/include/mal\_abm\_fast/config.hpp, mal-core/src/mal\_core/abm/include/mal\_abm\_fast/wire.hpp … (+2 more) \\
\code{ingest.coastline\_buffer\_m} & 2000.0 & m & mal-commonlib/src/mal\_commonlib/data/loaders/coastline.py, mal-core/src/mal\_core/abm/include/mal\_abm\_fast/config.hpp, mal-core/src/mal\_core/ingest/daily\_nc.py \\
\code{ingest.coastline\_resolution} & "f" & GSHHG & mal-commonlib/src/mal\_commonlib/data/loaders/coastline.py, mal-core/src/mal\_core/abm/include/mal\_abm\_fast/config.hpp \\
\code{ingest.fill\_frac\_valid\_max} & 0.5 & fracción & mal-core/src/mal\_core/abm/include/mal\_abm\_fast/config.hpp, mal-core/src/mal\_core/ingest/daily\_nc.py \\
\code{ingest.grid\_degrees\_step\_max} & 1 & ° & mal-core/src/mal\_core/abm/include/mal\_abm\_fast/config.hpp, mal-core/src/mal\_core/ingest/daily\_nc.py \\
\code{ingest.land\_mask\_threshold} & 0.5 & fracción & (loaded, never used) \\
\code{ingest.min\_dilation\_radius\_cells} & 1 & celdas & (loaded, never used) \\
\code{ingest.nc\_compression\_level} & 4 & nivel & mal-core/src/mal\_core/abm/include/mal\_abm\_fast/config.hpp, mal-core/src/mal\_core/ingest/daily\_nc.py \\
\code{ingest.nodata\_sentinel} & -9999.0 & float & mal-commonlib/src/mal\_commonlib/data/host\_utils.py, mal-commonlib/src/mal\_commonlib/data/loaders/\_legacy/hydrolakes.py … (+19 more) \\
\code{ingest.nodata\_sentinel\_int} & -9999 & int & mal-commonlib/src/mal\_commonlib/data/host\_utils.py, mal-commonlib/src/mal\_commonlib/data/loaders/dynamic\_world.py … (+2 more) \\
\code{ingest.ocean\_mask\_threshold} & 0.5 & fracción & (loaded, never used) \\
\code{ingest.residential\_from\_indoor\_frac} & 0.8 & fracción & mal-core/src/mal\_core/abm/include/mal\_abm\_fast/config.hpp, mal-core/src/mal\_core/abm/src/host\_landscape.cpp \\
\code{ingest.sea\_water\_floor} & 0.05 & fracción & mal-commonlib/src/mal\_commonlib/data/loaders/coastline.py, mal-core/src/mal\_core/abm/include/mal\_abm\_fast/config.hpp … (+2 more) \\
\code{ingest.valid\_count\_epsilon} & 1e-09 & cuentas & mal-commonlib/src/mal\_commonlib/data/loaders/dynamic\_world.py, mal-commonlib/src/mal\_commonlib/data/loaders/jrc\_gsw.py, mal-core/src/mal\_core/abm/include/mal\_abm\_fast/config.hpp \\
\code{pluvial.cold\_start\_z} & 0.0 & adimensional & mal-core/src/mal\_core/abm/include/mal\_abm\_fast/config.hpp, mal-core/src/mal\_core/abm/src/coordinator.cpp \\
\code{pluvial.delta\_days} & 1.0 & días & mal-core/src/mal\_core/abm/include/mal\_abm\_fast/config.hpp, mal-core/src/mal\_core/abm/include/mal\_abm\_fast/habitat\_availability.hpp, mal-core/src/mal\_core/abm/src/coordinator.cpp \\
\code{pluvial.dry\_memory\_days} & 5.0 & días & mal-core/src/mal\_core/abm/include/mal\_abm\_fast/config.hpp, mal-core/src/mal\_core/abm/include/mal\_abm\_fast/coordinator.hpp, mal-core/src/mal\_core/abm/include/mal\_abm\_fast/habitat\_availability.hpp \\
\code{pluvial.rain\_scale\_mm} & 2.5 & mm & mal-core/src/mal\_core/abm/include/mal\_abm\_fast/config.hpp, mal-core/src/mal\_core/abm/include/mal\_abm\_fast/habitat\_availability.hpp \\
\code{urban.beta\_urban} & 0.75 & adimensional & mal-core/src/mal\_core/abm/include/mal\_abm\_fast/config.hpp, mal-core/src/mal\_core/abm/include/mal\_abm\_fast/habitat\_availability.hpp \\
\code{urban.rain\_pools\_built\_damage} & 0.3 & adimensional & mal-core/src/mal\_core/abm/include/mal\_abm\_fast/config.hpp, mal-core/src/mal\_core/abm/include/mal\_abm\_fast/habitat\_availability.hpp \\
\code{wildlife.class\_suitability.bare} & 0.2 & fracción & mal-commonlib/src/mal\_commonlib/data/loaders/wildlife.py \\
\code{wildlife.class\_suitability.built} & 0.05 & fracción & mal-commonlib/src/mal\_commonlib/data/loaders/wildlife.py \\
\code{wildlife.class\_suitability.crops} & 0.3 & fracción & mal-commonlib/src/mal\_commonlib/data/loaders/wildlife.py \\
\code{wildlife.class\_suitability.flooded\_vegetation} & 0.45 & fracción & mal-commonlib/src/mal\_commonlib/data/loaders/wildlife.py \\
\code{wildlife.class\_suitability.grass} & 0.6 & fracción & mal-commonlib/src/mal\_commonlib/data/loaders/wildlife.py \\
\code{wildlife.class\_suitability.shrub\_and\_scrub} & 0.7 & fracción & mal-commonlib/src/mal\_commonlib/data/loaders/wildlife.py \\
\code{wildlife.class\_suitability.snow\_and\_ice} & 0.1 & fracción & mal-commonlib/src/mal\_commonlib/data/loaders/wildlife.py \\
\code{wildlife.class\_suitability.trees} & 0.9 & fracción & mal-commonlib/src/mal\_commonlib/data/loaders/wildlife.py \\
\code{wildlife.class\_suitability.water} & 0.0 & fracción & mal-commonlib/src/mal\_commonlib/data/loaders/wildlife.py \\
\code{wildlife.weights.habitat} & 0.5 & fracción & mal-commonlib/src/mal\_commonlib/data/loaders/wildlife.py \\
\code{wildlife.weights.remoteness} & 0.2 & fracción & mal-commonlib/src/mal\_commonlib/data/loaders/wildlife.py \\
\code{wildlife.weights.water} & 0.3 & fracción & mal-commonlib/src/mal\_commonlib/data/loaders/wildlife.py \\
\bottomrule
\end{longtable}
% Generated by tools/gen_concepts_values.py -- do not edit.
\paragraph{The loaders' native grids (the resolution each product is published at)}
\begin{longtable}{@{}llll@{}}
\toprule
Loader & Native grid & Product / epoch & Source \\
\midrule
\endhead
\code{buildings} & --- & --- & \code{buildings.py} \\
\code{chirps} & 0.05° (\textasciitilde{}5 km) & v2.0 / from 1981 & \code{chirps.py} \\
\code{coastline} & --- & 2.3.7 & \code{coastline.py} \\
\code{dynamic\_world} & native: 10 m / 100 m (request scale) & GOOGLE/DYNAMICWORLD/V1 & \code{dynamic\_world.py} \\
\code{era5} & 0.1 deg & derived-era5-land-daily-statistics & \code{era5.py} \\
\code{ghsl} & \textasciitilde{}1 km & GHS\_SMOD\_E2030\_GLOBE\_R2023A & \code{ghsl.py} \\
\code{glw} & \textasciitilde{}10 km, \textasciitilde{}5 arc-minutes & GLW4 2020 & \code{glw.py} \\
\code{jrc\_gsw} & 30 m & 1984-2021 & \code{jrc\_gsw.py} \\
\code{modis} & 1 km & v061 & \code{modis.py} \\
\code{smap} & 0.25° & SMAP\_RSS\_L3\_SSS\_SMI\_MONTHLY\_V6 & \code{smap.py} \\
\code{wildlife} & --- & --- & \code{wildlife.py} \\
\code{worldpop} & \textasciitilde{}100 m & Global1 / 2000–2020 & \code{worldpop.py} \\
\bottomrule
\end{longtable}
\noindent\emph{---}: the loader's code states no such value; the cell is left empty rather than typed by hand.


\smallskip \noindent The generated file carries two tables. The first is the parameters table
above; the second tabulates the native grid at which each loader publishes its product,
together with the product or epoch it comes from and the loader's source file. The input
inventory of \S\ref{sec:envinputs} takes its native resolutions from that second table rather
than restating them here. A loader whose code states no such resolution has an empty cell
rather than a hand-typed one.

\smallskip \noindent The parameters table is not a list of what this stage reads: its last
column, the
``Read by'' column, is the answer to that question. The column is scanned from the two access
patterns the configuration actually uses --- the Python \code{model\_constant("section",
"key")} call and the C++ loader's \code{L.key(node, "section", "key")} ---
\src{tools/gen\_concepts\_values.py:197--218}. A cell that reads \emph{(loaded, never used)}
means that the key reaches the engine's configuration struct and that nothing acts on it: there
are exactly four such constants in the whole configuration,
\code{ingest.land\_mask\_threshold}, \code{ingest.min\_dilation\_radius\_cells},
\code{ingest.ocean\_mask\_threshold} and \code{contracts.env\_nc}. For the three \code{ingest}
thresholds no production code reads them at all, so the values they carry are inert. For
\code{contracts.env\_nc} the verdict is about the engine side --- the engine's loader reads it
into the configuration struct
\src{mal-core/src/mal\_core/abm/src/config.cc:496--497} and no engine code acts on it --- while
the Python builder reads the same key through the constants dictionary and writes it into the
environment NetCDF's contract attribute
\src{mal-core/src/mal\_core/ingest/daily\_nc.py:53},
\src{mal-core/src/mal\_core/ingest/daily\_nc.py:438}; that access pattern is not one the
scanner matches. The remaining rows name constants that belong to the stage's namespaces but
are read elsewhere: the download stage's loaders (\code{wildlife.*},
\code{ingest.coastline\_resolution}, \code{ingest.valid\_count\_epsilon},
\code{ingest.nodata\_sentinel\_int}), the C++ \code{host\_landscape}
(\code{ingest.residential\_from\_indoor\_frac}), and the engine's configuration and geodesy
(\code{contracts.cogs\_state}, \code{geodesy.*}).

\medskip \noindent The output contracts, as generated from the contract definition:

% Generated by tools/gen_concepts_values.py -- do not edit.
\paragraph{The output contracts}
\begin{longtable}{@{}llllll@{}}
\toprule
Artefact & Producer & Consumers & Version & Bands / variables & Sentinel \\
\midrule
\endhead
\code{env} & \code{mal\_core.ingest.env.build\_env\_tensor} & abm, prediction, scoring & "3.0" & rainfall, air\_temp\_c, jrc\_persistent, lc\_built, lc\_wetland, lc\_open\_plus\_water, lc\_crop, land\_frac (optional), salinity\_ppt (optional) & none \\
\code{habitat} & \code{mal\_core.ingest.env.build\_env\_tensor} & abm &  & hab\_type, habitat\_potential, row, col & none \\
\code{state\_cog} & \code{mal\_abm\_fast::write\_state\_cog} & scoring, prediction, training & "4.0" & adult\_patch\_count, adult\_cell\_count & -9999.0 \\
\code{state\_sidecar} & \code{mal\_abm\_fast::write\_state\_sidecar} & scoring & "4.0" &  &  \\
\code{manifest} & \code{mal\_core.download.manifest.update\_dataset} & download, ingest, abm &  &  &  \\
\code{host\_static} & \code{mal\_core.ingest.hosts.build\_host\_dataset} & abm &  & humans, livestock, urban, buildings &  \\
\code{mobility\_day} & \code{mal\_core.ingest.mobility.build\_mobility\_dataset} & abm &  &  &  \\
\bottomrule
\end{longtable}

\paragraph{The state sidecar's keys, in emission order}
\begin{longtable}{@{}l@{}}
\toprule
Key \\
\midrule
\endhead
\code{crs} \\
\code{transform} \\
\code{aoi\_slug} \\
\code{scale} \\
\code{year} \\
\code{month} \\
\code{seed} \\
\code{n\_rollouts} \\
\code{rollout\_index} \\
\code{generator\_version} \\
\code{abm\_params\_hash} \\
\code{contract\_version} \\
\code{band\_names} \\
\code{nodata} \\
\code{shape} \\
\bottomrule
\end{longtable}


\section{What is not answered, and what is not implemented} \label{sec:open}

The following are open, and are recorded here rather than left as silent omissions. Each is a
question a reader is entitled to ask and the code does not settle.

\begin{enumerate}
\item \textbf{The host path downloads; the environment path does not.} Is the asymmetry
intended? (\S\ref{sec:hostdownloads}.) 
\item \textbf{There are two different totals of the
habitat potential} --- the Python one and the engine's. Which is the intended definition, and
which should the GeoPackage's amplitude use? (\S\ref{sec:habitat}.) 
\item \textbf{The GeoPackage
carries no version.} Should a contract version be stamped on it? (\S\ref{sec:gpkg}.) 
\item
\textbf{The \code{dynamics} column is written as a string and read as an integer}, so the
engine's per-patch flag is always zero. (\S\ref{sec:gpkg}.) 
\item \textbf{The manifest's habitat
variable list does not match the file}. (\S\ref{sec:manifest}.) 
\item \textbf{The command-line
defaults place files and manifest entries in different directories}. (\S\ref{sec:entry}.) 
\item
\textbf{The environment NetCDF is registered under one year key although it spans several
years}. (\S\ref{sec:manifest}.) 
\item \textbf{Some numbers are code literals and not
constants-file keys}, and are generated by the code-constant family rather than by the
constants file. (\S\ref{sec:literals}.) 
\item \textbf{There
are two capacity scales in the codebase}, the engine's and the Python diagnostic one.
(\S\ref{sec:dynamic}.) 
\item \textbf{The three candidate floors have distinct roles}, but
nothing prevents the GeoPackage from carrying patches the engine immediately drops.
(\S\ref{sec:habitat}.) 
\item \textbf{The open-sea filter degrades silently} when the coastline
artefact is missing. (\S\ref{sec:statics}.) 
\item \textbf{The declared \code{host\_static} slot
is never read} by the environment builder. (\S\ref{sec:envinputs}.) 
\item \textbf{The frame has
two owners}, reconciled only by a shape assertion on the GeoPackage path. (\S\ref{sec:frame}.)
\item \textbf{The cell-area factor is an approximation} whose size is stated rather than
corrected. (\S\ref{sec:frame}.) 
\item \textbf{The wind product is not produced here}; it comes
from the download stage. (\S\ref{sec:seam}.) 
\item \textbf{Whether the static planes should
follow the run year} is not settled: multi-year land cover is averaged. (\S\ref{sec:statics}.)
\item \textbf{The urban-class reclassification is inverted} with respect to the GHS-SMOD product
it reclassifies, so the engine reads the area of interest as mostly urban when it is mostly
rural. (\S\ref{sec:hosts}.)
\item \textbf{The delivered environment NetCDF was produced by the pre-fix implementation},
which is why it carries no \code{land\_frac} plane. (\S\ref{sec:envnc}.)
\item \textbf{The delivered GeoPackage's candidate set includes open-sea cells} that the
current open-sea rule would remove. (\S\ref{sec:gpkg}.)
\end{enumerate}

\section{Evidence appendix} \label{sec:appendix}

\subsection{Implementation notes a replicator needs}

\begin{enumerate}
\item \textbf{The frame is set by the rainfall file}, so a replicator who changes the AOI must
rebuild that file on the new grid (\S\ref{sec:frame}). 
\item \textbf{The environment NetCDF is
written before the GeoPackage}, and the GeoPackage is derived from the file that was written, so
the candidate set can only be reproduced by reproducing the statics first (\S\ref{sec:order}).
\item \textbf{The NoData fill is spatial, not temporal}: changing it means changing the
algorithm at \src{mal-core/src/mal\_core/ingest/daily\_nc.py:206--223}, not a parameter. 
\item
\textbf{Open sea is a distance query, not a dilation}, and the distance is computed from the
land-fraction plane \src{mal-commonlib/src/mal\_commonlib/data/loaders/coastline.py:246--261}.
\item \textbf{The coastal band can be overridden per run} through the
\code{COASTLINE\_BUFFER\_M} environment variable; a negative value disables the filter
\src{mal-core/src/mal\_core/ingest/daily\_nc.py:240--250}. 
\item \textbf{The temperature channel
is re-indexed by day, not by position} (\S\ref{sec:channels}). 
\item \textbf{Static planes are
averaged over a year axis} when the source file has one (\S\ref{sec:statics}). 
\item \textbf{The
GeoPackage's point geometry is the cell centre}, not the cell polygon (\S\ref{sec:gpkg}). 
\item \textbf{The engine does
not need the \code{hab\_<type>} columns}, so removing them from the writer would not change the
simulation (\S\ref{sec:gpkg}).
\end{enumerate}

\subsection{Values that cannot be generated} \label{sec:literals}

The values below are literals in the source, or belong to a provider's own specification, and
they have no key either in the shared constants file or in the generated code-constant table.
They are quoted in this document as readings of the code or of the provider, and they are the
ones a reviewer should check against the cited lines. Every other number this document quotes
is a generated macro.

\begin{center}
\small
\begin{longtable}{@{}p{4.6cm}p{2.6cm}p{7.8cm}@{}}
\toprule
Value & Reading & Where \\
\midrule
\endhead
Generator version & \code{proxy-habitat-v1-0.1.0}
  & \src{mal-core/src/mal\_core/ingest/daily\_nc.py:439} \\
GHS-SMOD urban class scheme & codes $10$--$30$
  & \src{mal-commonlib/src/mal\_commonlib/data/loaders/ghsl.py:76--79} \\
Engine indoor-fraction threshold & $40$
  & \src{mal-core/src/mal\_core/abm/src/host\_landscape.cpp:111--118} \\
WorldPop product year & $2019$
  & \src{mal-core/src/mal\_core/ingest/hosts.py:40} \\
\bottomrule
\end{longtable}
\end{center}

\subsection{What is verified, and what is not}

\emph{Verified against the delivered artefacts} (read directly from the files in
\code{data/ghana/}):

\begin{itemize}
\item the environment NetCDF's variables, shapes, dtypes, global attributes, fill value and
absence of any coordinate reference system; 
\item the GeoPackage's feature count, column set,
column types, the dominant-type and dynamics distributions and the amplitude range; 
\item the
manifest entries for the environment NetCDF, the habitat GeoPackage, the host grid,
including the mismatch between the habitat entry's declared variables and
the file's actual columns, and the absence of the required-for-engine flag on the first two;
\item the input files' own variable lists and attributes, including the absence of coordinate
variables in the rainfall and temperature files and the Dynamic World class mapping recorded in
the file's attributes.
\end{itemize}

\emph{Verified against the source}: every behavioural claim in this document, by reading the
cited lines.

\emph{Not verified, and not claimed}:

\begin{itemize}
\item that the numbers in the delivered artefacts were produced by the current code --- they were
produced by the pre-fix implementation, as \S\ref{sec:gpkg} and \S\ref{sec:envnc} state; 
\item the exact magnitude of the
cell-area bias, which was measured outside the codebase for the Ghana box during the
preparation of this document and is reported in \S\ref{sec:frame} as an order of magnitude
rather than as a generated constant;
\item the downstream behaviour of the engine on the delivered files, which is a question for the
simulation stage's document; 
\item whether the delivered manifest entries were produced by the
registration calls in this document, given the required-for-engine flag discrepancy noted in
\S\ref{sec:entry}.
\end{itemize}

\subsection{What the artefacts must not be used for}

\begin{itemize} 
\item The habitat potential must not be read as a pool count, a depth, a volume
or a probability (\S\ref{sec:hnot}). 
\item The livestock planes must not be read as head counts,
and the building fraction must not be read as an area fraction (\S\ref{sec:hosts}). 
\item The
candidate count printed by the ingest build must not be read as the engine's candidate count
(\S\ref{sec:habitat}). 
\item The \code{dynamics} column of the GeoPackage must not be read as a
faithful round trip into the engine (\S\ref{sec:gpkg}). 
\item The environment NetCDF's
coordinates must not be read as geographic coordinates without checking the file
(\S\ref{sec:frame}).
\end{itemize}

\end{document}
